\documentclass[11pt,a4paper]{ltjsarticle}

% 日本語査読用：LuaLaTeX + LuaTeX-ja
% 英語版 main.tex (v72) からの翻訳。数式・label・cite は英語版と共通に保つ。
\usepackage{amsmath,amssymb,amsthm,mathtools}
\numberwithin{equation}{section}
\usepackage{enumitem}
\usepackage{array}
\usepackage{tabularx}
\usepackage{geometry}
\geometry{margin=1in}
\usepackage{cite}
\usepackage{hyperref}
\usepackage{url}
\usepackage{tikz}
\hypersetup{colorlinks=true,linkcolor=blue,citecolor=blue,urlcolor=blue}
\setlength{\emergencystretch}{3em}

\theoremstyle{plain}
\newtheorem{theorem}{定理}[section]
\newtheorem{lemma}[theorem]{補題}
\newtheorem{proposition}[theorem]{命題}
\newtheorem{corollary}[theorem]{系}
\theoremstyle{definition}
\newtheorem{definition}[theorem]{定義}
\theoremstyle{remark}
\newtheorem{remark}[theorem]{備考}
\newtheorem{example}[theorem]{例}
\newtheorem{openproblem}[theorem]{未解決問題}
\renewcommand{\proofname}{証明}

\newcommand{\Wit}{\mathsf{W}}
\newcommand{\Cmcf}[2]{\mathcal{C}^{\mathrm{mcf}}_{#1,#2}}
\newcommand{\D}{\mathcal{D}}
\newcommand{\out}{\mathrm{out}}
\newcommand{\tuple}[1]{\vec{#1}}
\newcommand{\blank}{\square}
\newcommand{\eps}{\varepsilon}
\newcommand{\pausesym}{\bot}
\newcommand{\fan}{\operatorname{fan}}
\newcommand{\YSub}[1]{\mathsf{YSub}(#1)}
\newcommand{\YSL}[2]{\mathsf{YSL}(#1,#2)}
\newcommand{\evalmap}{\mathsf{ev}}

\newcommand{\cutpos}[1]{%
  \raisebox{-0.7ex}{%
    \tikz[baseline=(char.base)]{%
      \node[draw,circle,inner sep=0.4pt,font=\scriptsize]
        (char) {\ensuremath{#1}};%
    }%
  }%
}

\title{固定有限モノイド型付けの下での多重文脈自由言語の正例学習}
\author{Takayuki Kuriyama\\
独立研究者（東京）\\
\texttt{growup.kuriyama@gmail.com}}
\date{}

\begin{document}
\maketitle

\begin{abstract}
有限モノイドへの準同型 $h:\Sigma^*\to M$ とファンアウト上界 $f$ を
固定する．対応する成分の $h$-値が等しい場合に限りタプルを交換できる
多重文脈自由言語の正例学習を研究する．規則ランクを制限すると，
このクラスは Yoshinaka の代入可能な signature に関する一般理論の一例となる．
本稿では，型付けで制限した観測タプルからの文法構成との対応を検証し，
ランク上界を設けない構成と Yoshinaka の学習器との関係も明らかにする．

本稿の中心的な定量的問題は，元の MCFG の文法サイズと規則厚みで測ったときの，
正例特徴データの\emph{総文字数}である．固定接頭辞--接尾辞窓で型が決まる場合，
ファンアウトと規則ランクを固定すれば，両パラメータに関する多項式上界が得られる．
その証明では，タプルの複数の成分にわたる境界情報を同時に保存する．
さらに，有限アーベル群への準同型値と固定窓を組み合わせた型付けにも，
零和による短縮を用いて多項式上界を与える．
規則ランクが1なら，任意の固定有限型付けについて，
文法サイズのみの多項式で特徴データを抑えられる．
他方，任意の有限型付けに対して元文法の厚みによる多項式上界が
一様に成立するわけではないことは，すでに CFG で知られている．
本稿の窓型付けと窓・アーベル群型付けの条件は，この障害を回避するものである．
また，有限型付け，正則フィルタおよびファンアウトに関する分離を示し，
最小ファンアウトと最小規則ランクがともに2である，
境界型付けされた非文脈自由標的を構成する．
\end{abstract}

\medskip
\noindent\textbf{キーワード：}
文法推論；多重文脈自由文法；正例データ；有限モノイド型付け；
多次元代入可能性；Gold の極限同定；規則ランク
\medskip

\section{序論}
\label{sec:introduction}

正例からの文法推論は，言語の正例だけを列挙するデータから，
いつその言語を復元できるかを問う~\cite{Gold1967,Angluin1980}．
分布的手法は，部分文字列やタプルについて観測された交換可能性を
文法の非終端記号へと変換することで，制限された文法クラスの学習を可能にする．
文脈自由文法（CFG）の場合，これには代入可能言語~\cite{ClarkEyraud2007} と
Yoshinaka の $(k,\ell)$ 階層~\cite{Yoshinaka2008} が含まれる．
Yoshinaka は同じ着想を多重文脈自由文法（MCFG）へ拡張した．
この場合，非終端記号はタプルを生成し，代入は複数の穴をもつ
文脈によって検査される~\cite{Yoshinaka2009,Yoshinaka2011}．
MCFG は代表的な弱文脈依存形式体系であり，
不連続依存や交差直列依存を直接表現できる
~\cite{SekiEtAl1991,Kallmeyer2010}．

\paragraph{一般学習原理とその適用範囲．}
Yoshinaka の2011年の観測タプルに基づく学習器が，
型付けを導入しない MCFG 学習の基準となる~\cite{Yoshinaka2011}．
その後の signature に関する一般理論は，
アリティの制限された MCFG，ソートの有限精密化，接頭辞--接尾辞型付けを含み，
すでに正例からの同定，多項式時間更新，および個数が有界な特徴集合の存在を与える
~\cite[Theorem~1, Lemma~3 and Examples~3--4]{Yoshinaka2015General}．
規則ランクが有界な場合，本稿の付録~\ref{sec:general-framework} は，
型付き言語クラスと signature の間の厳密な特徴づけと，
再構成作用素同士の対応を証明する．
したがって，\emph{質的な}学習原理そのものは先行研究である．
本稿の問題は，元の MCFG 表示に相対的な正例特徴データの
\emph{総文字数}，および有限型付け・ファンアウト・規則ランクに関する分離である．

本稿では学習対象クラスの一部として，有限モノイドへの準同型
$h:\Sigma^*\to M$ を固定する．観測されたタプルは，
対応する成分の $h$-値が等しい場合に限って比較する．
合成可能な型による制約は正例テキストが与えられる前から固定されており，
有限なバイアスの選択と，そのバイアスの下での学習を分離する
~\cite{CosteTyping2004}．
重要な二つの例は接頭辞--接尾辞準同型と DFA の遷移準同型である．
固定接頭辞--接尾辞窓は，個々の標本に応じて調整するパラメータではなく，
標的クラスに対して独立かつ一様に選んだ有限パラメータである．
本稿の総文字数に関する主上界は，ファンアウトと規則ランクを指定すれば，
その窓型付けクラスのすべての標的に適用される．
ファンアウトもバイアスの一部である．これは代入可能性を検査する
タプルのアリティを決定するからである．
したがって，固定型付けは有限状態の情報を供給し，
非有界な依存関係は MCFG が担う．

他の学習バイアスは異なる選択を行う．
Takada は導出に対する正則制御言語を用いる
~\cite{Takada1988,Takada1995}．
一方，局所的代入可能性は因子を比較するための証拠を制限する
~\cite{CosteGaretNicolas2012,CosteNicolas2019}．
正例からの学習は他の弱文脈依存形式体系についても研究されている
~\cite{OatesEtAl2006}．
より豊富な情報を利用する MCFG 学習には，
所属質問や双対文脈集合を用いる手法がある
~\cite{Yoshinaka2010,ClarkYoshinaka2014,KanazawaYoshinaka2023}．
本稿はこれらと異なり正例データのみに限定し，
観測タプルの値を標本由来の基本対象とし，
固定有限型付けにより代入を制限する．

\paragraph{本稿の貢献．}
本稿が扱うのは，有限特徴集合の存在やその個数だけではなく，
正例特徴データの\emph{総文字数}である．
\begin{enumerate}[label=\textup{(\arabic*)},leftmargin=2.4em]
\item \emph{複数成分に関する定量的上界．}
固定したファンアウトと規則ランクの下で，
固定接頭辞--接尾辞窓により型が決定されるとき，
縮約済み良形表示 $G$ に対し，特徴データの総文字数は
$O(\|G\|^2\max\{1,\theta_G\})$ 以下となる
（定理~\ref{thm:boundary-polydata}）．
鍵となるのは，タプルの全成分にわたる境界情報を同時に保存する
短いタプルの補題である．
付録~\ref{app:window-abelian} は，固定有限アーベル群の値を窓と組み合わせた場合へ
結果を拡張し，総文字数
$O(\|G\|^3\bar\theta_G)$ を与える
（定理~\ref{thm:window-abelian-polydata}）．

\item \emph{型付け・ファンアウト・規則ランクに関する分離．}
Yoshinaka の型なしクラスに対する任意の正則フィルタ断片の外にある言語
（定理~\ref{thm:arbitrary-filter-separation}），
ファンアウト1では型付き代入可能だが，
ファンアウト2では\emph{いかなる}有限型付けでも代入可能でない標的
（命題~\ref{prop:strong-fanout-separation}），
および最小ファンアウト・最小規則ランクがともに2である
境界型付け非文脈自由言語
（定理~\ref{thm:intrinsic-ranktwo-witness}）を与える．
\end{enumerate}
学習作用素そのものは，Yoshinaka の観測タプル構成を
成分ごとに型付けしたものであり，
規則ランク有界の場合に一般 signature 枠組みに含まれることは
命題~\ref{prop:yoshinaka2015-embedding} で示す．

\paragraph{先行研究との関係と本稿の範囲．}
観測タプル文法，保守的な更新方針およびランク上界を設けない質的同定は
Yoshinaka~\cite{Yoshinaka2009,Yoshinaka2011} による．
一般的な有限ソートの枠組みと特徴集合の個数に関する上界は
Yoshinaka~\cite{Yoshinaka2015General} の成果である．
本稿の直接的な健全性・完全性の証明は，
成分ごとに型付けした特定の文法構成を検証するものであって，
新しい一般的学習機構を主張するものではない．
規則ランク1における特徴データ長の結果は基本的な精密化の帰結である．
より本質的な高ランクの成果は，複数のタプル成分について
境界を同時に保存することで得られる上界
（定理~\ref{thm:boundary-polydata}）である．

ファンアウト1の固定接頭辞--接尾辞型付けに関する多項式特徴データは，
Yoshinaka~\cite[Theorem~1]{Yoshinaka2008} によって得られている．
関連する CFG 論文~\cite{kuriyamaCFG} は任意の有限型付けを扱い，
ファンアウト1で作用素に依存しない下界を証明する．
これらは先行結果であり，本稿独自の MCFG 定理としては掲げない．
本稿固有の問題は，タプルの成分ごとの型付け，
高ランクでの特徴データ長，およびファンアウト・正則フィルタとの相互作用である．
内在的ランクを示す言語の証明で利用する非分岐性の下界も先行研究に基づく．
本稿の寄与は，その境界型付けによる構成，
型付きクラスへの所属と Yoshinaka のクラスからの分離，
ならびに定量的定理との結び付きである．

多項式上界の各主張では，ファンアウト・規則ランク上界・有限型付けを固定する．
集合駆動型再構成の特徴データは\emph{言語レベル}の厳密性を保証し，
構文レベルの安定化は保守的な逐次学習器が担う．
型付き完全精密化と空成分除去に関する障害は，
それら特定の構成の障害であって，すべての学習器に対する下界ではない．
ファンアウト1における独立した作用素非依存の下界は，
関連 CFG 論文~\cite{kuriyamaCFG} の成果である．

第~\ref{sec:prelim}節から第~\ref{sec:char-data}節で，
モデル，再構成，収束および定量的上界を扱う．
第~\ref{sec:comparison}節では Yoshinaka の言語クラスとの比較と
分離結果を与える．

\paragraph{証明の方法とAI支援研究．}
本稿は理論研究であり，主張の根拠はここに提示する数学的証明であって，
計算機実験や自動化された査読ではない．
生成AIは候補となる論証の探索・批判的検討や原稿改訂の補助に利用したが，
その出力自体を数学的検証とみなすことはしない．
使用したツールとその役割は論文末尾で開示する．

\section{準備}
\label{sec:prelim}

以下，$\Sigma$ は有限の終端アルファベット，
$\eps$ は空語，$\Sigma^*$ は $\Sigma$ 上の有限語全体，
$\Sigma^+:=\Sigma^*\setminus\{\eps\}$ とする．

\subsection{正例データからの学習}

Gold の説明的学習モデルを採用し，MCFG に関する議論を自己完結にするため，
必要な規約を再掲する．

\begin{definition}[テキストと TxtEx 同定]
Gold の説明的同定の基準を用いる~\cite{Gold1967}．
空言語を扱うため，
$\pausesym\notin\Sigma$ なる停止記号を固定する．
したがって $\pausesym,\pausesym,\ldots$ は空言語のテキストである．
$L\subseteq\Sigma^*$ の\emph{テキスト}とは，
$L\cup\{\pausesym\}$ 上の無限列で，
$L$ のすべての要素が少なくとも一度出現するものをいう．
停止記号はいかなる位置にも出現してよく，
$\pausesym$ はテキストの終了を示す記号ではない．

有限接頭部分の\emph{内容}とは，そこに出現した
$\Sigma^*$ 上の正例からなる有限集合 $K\subseteq L$ である．
したがって，停止記号の出現位置や個数，
同じ正例の繰返し回数は内容に影響しない．
学習器は各有限接頭部分から仮説文法への計算可能関数である．

学習器が $L$ を正例データから\emph{TxtEx 同定}するとは，
$L$ の任意のテキストに対して，ある段階 $n_0$ と
一つの仮説文法 $H$ が存在し，
$n_0$ 以後のすべての段階で出力が文字どおり $H$ に等しく，
かつ $L(H)=L$ となることをいう．
クラスのすべての言語をこの意味で同定するとき，
そのクラスを TxtEx 同定するという．
逐次学習器が\emph{保守的}であるとは，
新しく観測された正例が現在の仮説言語に既に含まれているなら
仮説を維持し，生成されていないときに限って仮説を変更することをいう．
\end{definition}

以下，$L(H_0)=\emptyset$ を満たす
特別な仮説コード $H_0$ を固定する．

\begin{definition}[集合駆動型再構成作用素と特徴標本]
\label{def:characteristic-sample}
\emph{集合駆動型再構成作用素} $\mathcal B$ は，
各有限正例標本 $K$ を文法 $\mathcal B(K)$ に写し，
標本の提示順序や重複度ではなく，集合 $K$ だけに依存する．
すべての有限 $K$ に対して
$K\subseteq L(\mathcal B(K))$ が成り立つとき，
$\mathcal B$ は\emph{標本整合的}であるという．

de~la~Higuera~\cite{deLaHiguera1997} の
特徴集合の見方に従い，
標的言語 $L$ の有限集合 $K^\star\subseteq L$ が
$\mathcal B$ の\emph{特徴標本}であるとは，
任意の有限集合 $K$ に対し
$K^\star\subseteq K\subseteq L$ ならば
$L(\mathcal B(K))=L$ が成立することをいう．
$K^\star=\emptyset$ の場合も許す．
\emph{特徴データ}とは，このような標本の正例内容の略称として用いる．
多項式上界を主張する際には，
集合の要素数を抑えるのか，正例の総文字数を抑えるのかを明記する．

この有限標本の概念は TxtEx 学習と異なる．
蓄積標本が増加するとき，
$\mathcal B(K)$ が構文的に安定する必要はない．
逐次学習器は，厳密な再構成を証明した後で別途定義する．
\end{definition}

\begin{definition}[正例標本のサイズ]
有限集合 $K\subseteq\Sigma^*$ に対し，
\[
 \|K\|_+:=\sum_{w\in K}(|w|+1)
\]
と定義する．
タプル $\tuple u=(u_1,\ldots,u_d)$ の成分総長を
\[
 \|\tuple u\|_1:=\sum_{i=1}^d |u_i|
\]
と書く．
\end{definition}

\subsection{有限モノイド型付け}

\begin{definition}[有限型付けの表示]
\label{def-monoid-morphism}
有限モノイド準同型 $h:\Sigma^*\to M$ は，
$M$ の乗法表と各文字の像 $h(a)$（$a\in\Sigma$）によって表示する．
したがって $h(w)$ は左から右へ積を取ることで，
$|w|$ に線形な時間で計算できる．
タプル $\tuple w=(w_1,\ldots,w_d)$ に対して
\[
 h^{(d)}(\tuple w):=(h(w_1),\ldots,h(w_d))\in M^d
\]
と書く．
\end{definition}

\begin{definition}[有限型付けの精密化]
\label{def:h-refinement}
$h:\Sigma^*\to M$ と $h':\Sigma^*\to M'$ を
有限モノイド準同型とする．
\[
  \ker h'\subseteq\ker h
\]
が成り立つとき，$h'$ は $h$ を\emph{精密化する}といい，
$h\preceq h'$ と書く．
同値な条件として，値域をそれぞれの像に制限すると，
一意なモノイド準同型
$\upsilon:\operatorname{im}h'\to\operatorname{im}h$
が存在して，
$h=\upsilon\circ h'$ となる．
\end{definition}

二つの型付け $h,g$ の積型付け $(h,g)$ は両方を精密化し，
$\ker(h,g)=\ker h\cap\ker g$ を満たす．

ここでは有限モノイド準同型を有限状態型付けとして用いる．
$h_{\mathrm{triv}}:\Sigma^*\to\{1\}$ を自明型付けとする．
標的を含む正則言語の DFA の遷移準同型も例である
~\cite{Pin1986}．
利用するのは準同型による有限な要約だけであり，
所属質問へのアクセスではない．
$h^{(d)}(\tuple w)$ を\emph{生成語型}と呼ぶ．
以下に定義する接頭辞--接尾辞準同型は，
後で使う境界型付けを与える．

\begin{definition}[接頭辞--接尾辞型付け]
\label{def:prefix-suffix-typing}
$k,\ell\ge0$ に対し $m_0:=\max\{1,k+\ell\}$ とおく．
$|w|\ge m_0$ の語 $w$ について，
$\operatorname{pref}_k(w)$ と $\operatorname{suf}_\ell(w)$ は，
それぞれ指定した長さの接頭辞と接尾辞であり，
窓の長さが0なら $\eps$ とする．
\[
 h_\Sigma^{k,\ell}(w)=
 \begin{cases}
   w, & |w|<m_0,\\
   \langle\operatorname{pref}_k(w),\operatorname{suf}_\ell(w)\rangle,
      & |w|\ge m_0
 \end{cases}
\]
と定義し，短い語と接頭辞--接尾辞の対は
相異なるタグで区別する．
$M_{k,\ell}:=h_\Sigma^{k,\ell}(\Sigma^*)$ とおく．
短い語 $x,y$ と長い要素 $\langle u,v\rangle$ の積を
\[
\begin{aligned}
 x\cdot y&:=h_\Sigma^{k,\ell}(xy), &
 x\cdot\langle u,v\rangle&:=\langle\operatorname{pref}_k(xu),v\rangle,\\
 \langle u,v\rangle\cdot y&:=\langle u,\operatorname{suf}_\ell(vy)\rangle, &
 \langle u,v\rangle\cdot\langle u',v'\rangle&:=\langle u,v'\rangle
\end{aligned}
\]
によって定める．
\end{definition}

この積により
$h_\Sigma^{k,\ell}(xy)
=h_\Sigma^{k,\ell}(x)\cdot h_\Sigma^{k,\ell}(y)$
が成り立つ．
$h_\Sigma^{k,\ell}$ は全射なので，
積の結合性は連接から従い，
$h_\Sigma^{k,\ell}(\eps)$ は単位元である．
よって
$h_\Sigma^{k,\ell}:\Sigma^*\to M_{k,\ell}$
は有限モノイド準同型である．

有限型付け $h$ が
\emph{$(k,\ell)$-境界決定的}であるとは，
$h\preceq h_\Sigma^{k,\ell}$ が成り立つことをいう．
同値な条件として，
接頭辞--接尾辞型が等しければ $h$-型も等しく，
接頭辞--接尾辞型付けの像上に因子写像が存在する．
以下ではこれを単に\emph{境界型付け}とも呼ぶ．

\subsection{多重文脈自由文法}

多重文脈自由文法の標準的なタプル生成による定義を用いる
~\cite{SekiEtAl1991}．
MCFG と関連する弱文脈依存形式体系の概説は
\cite{Kallmeyer2010} を参照されたい．
以下の任意規則ランクに対応した記法を，
再構成と証明の全体で用いる．

\begin{definition}[標準的な MCFG 表示]
\label{def:standard-mcfg}
\emph{多重文脈自由文法}は
$G=(V,\Sigma,P,S,\fan)$ という組である．
ここで $V$ は有限非終端記号集合，
$\Sigma$ は有限終端アルファベット，
$S\in V$ は $\fan(S)=1$ を満たす開始記号，
$\fan(A)\ge1$ は各 $A\in V$ のファンアウトである．
ランク $r\ge0$ の規則 $R$ は
\[
  R:\quad A\to\rho(B_1,\ldots,B_r),
  \qquad
  \rho=(\alpha_1,\ldots,\alpha_{d_0})
\]
という形をもつ．ここで $\rho$ は出力テンプレート，
$d_0=\fan(A)$，$d_j=\fan(B_j)$ である．
\[
  \Xi_j:=\{\xi_j^1,\ldots,\xi_j^{d_j}\},
  \qquad
  \Gamma_\rho:=\Sigma\uplus\Xi_1\uplus\cdots\uplus\Xi_r
\]
とすると，各 $\alpha_i\in\Gamma_\rho^*$ であり，
全タプル $(\alpha_1,\ldots,\alpha_{d_0})$ において
$\Xi_1\uplus\cdots\uplus\Xi_r$ の変数は
それぞれ高々一度しか出現しない．
この高々一度という条件を\emph{線形性}という．
すべての変数がちょうど一度出現するとき，規則は\emph{非消去}である．
非消去性は生成タプルの各成分が非空であることを意味しない．
空語の成分は引き続き許す．

子タプル
$\tuple u_j=(u_{j,1},\ldots,u_{j,d_j})
\in(\Sigma^*)^{d_j}$ に対し，
各 $\xi_j^\kappa$ を $u_{j,\kappa}$ で同時に置き換えることで
$\rho(\tuple u_1,\ldots,\tuple u_r)\in(\Sigma^*)^{d_0}$
を定義する．
$r=0$ の場合，規則は定数規則であり，
タプル $(\alpha_1,\ldots,\alpha_{d_0})$ を生成する．
各タプル言語 $L_A(G)$ は，
すべての規則適用で閉じた，成分ごとの包含順序で最小の言語族である．
\[
 L(G):=\{w\in\Sigma^*:(w)\in L_S(G)\}.
\]
同値に，$A$ を根とする有限導出木を
規則テンプレートによって葉から評価する．
その生成結果が
$\tuple u\in(\Sigma^*)^{\fan(A)}$ となるとき
$A\Rightarrow_G^*\tuple u$ と書く．
これを\emph{終端導出}，$\tuple u$ をその\emph{終端タプル}という．
開始記号からの終端導出
$S\Rightarrow_G^*(w)$ を\emph{成功導出}と呼ぶ．
すべての非終端記号について $\fan(A)\le f$ のとき
$f$-MCFG といい，その生成言語を $f$-MCFL という．
すべての規則ランクが $r_{\max}$ 以下の文法は
\emph{規則ランク上界} $r_{\max}$ をもつ．
\end{definition}

ファンアウト1の場合，この形式体系には通常の文脈自由文法が含まれる．
\emph{線形文脈自由文法}とは，
各生成規則の右辺に高々一つの非終端記号しか出現しない CFG であり，
その生成言語を\emph{線形文脈自由言語}という．
これは，上で述べた MCFG テンプレートの「線形性」とは異なる．
MCFG の線形性は，各規則変数が高々一度しか出現しない条件である．

有限 MCFG の表示には，記号単位の明示的な標準符号化を固定し，
その長さを $\|G\|$ と書く．
符号化は，非終端記号・規則の記録と，
各規則テンプレート中に出現するすべての終端記号・変数を数える．
以下の文法サイズに関する多項式上界はすべてこの符号化に関するものである．

整数 $f\ge1$，$r\ge0$ に対し，
\[
\begin{aligned}
  \mathsf{MCFL}(f,r):=\{L(G):{}&
   G\text{ はファンアウトが高々 }f\text{ の MCFG,}\\
 &\text{かつ規則ランクが高々 }r\}
\end{aligned}
\]
と書く．
任意の有限 MCFG は有限の規則ランクをもつため，
$f\text{-}\mathrm{MCFL}
=\bigcup_{r<\infty}\mathsf{MCFL}(f,r)$ である．
後で「ファンアウト $f$ における最小規則ランク」と言うときは，
与えられた特定の文法のランクではなく，
ファンアウト $f$ のすべての表示の間での言語レベルの最小値を意味する．

以下，$f$ はファンアウト（次元）を，
$r$ は一般の規則ランクを，
$r_{\max}$ は事前指定したランク上界を表す．
規則テンプレートの変数は常に $\xi_j^\kappa$ と書き，
$x,y,z$（矢印付きも含む）は語や観測タプルに用いる．
例に現れる数字や $\mathtt\#$ を含む終端記号は
等幅フォントで表す．

ファンアウト $d$ が同じ非終端記号 $A,B$ に対して，
$A\to B$ は単項の恒等テンプレート規則
$A\to(\xi_1^1,\ldots,\xi_1^d)(B)$ の略記である．
これは成分間の恒等対応だけを表す．
例えばファンアウト2では
$A\to(\xi_1^1,\xi_1^2)(B)$ を意味し，
$A\to(\xi_1^2,\xi_1^1)(B)$ のような交換規則ではない．
同様に，$A\to\tuple c$ はランク0の定数規則を表す．
これらの略記を，型付き開始リンクと，
後述する再構成の Rule~(Link) に用いる．

標準形式ではランク0の規則，恒等でない単項規則，
任意の有限ランクの規則および消去規則を許す．
非消去規則
$R:A\to\rho(B_1,\ldots,B_r)$，
$\rho=(\alpha_1,\ldots,\alpha_{d_0})$ が
\emph{非置換}であるとは，
親の成分を左から右へ走査したとき，各子 $B_j$ の変数
$\xi_j^1,\ldots,\xi_j^{\fan(B_j)}$
がこの順序で出現することをいう．
Yoshinaka~\cite[Section~2.3]{Yoshinaka2011} に従い，
どの $\alpha_i$ も空でない規則を
\emph{$\lambda$-自由}と呼ぶ．
非置換規則のうち，同一の子の隣接変数
$\xi_j^\kappa\xi_j^{\kappa+1}$ が
どの $\alpha_i$ にも隣接して現れないものを
\emph{非融合}と呼ぶ．
すべての規則が $\lambda$-自由，非消去，非置換，非融合ならば，
表示は\emph{良形}である．
空語は一貫して $\eps$ と書き，
「$\lambda$-自由」は Yoshinaka の標準的な規則条件の名称としてだけ用いる．

完全性の証明では良形表示を用いる．
非融合条件により，ある子の兄弟部分の生成語を固定しても，
その子の隣接する成分は，非空の終端部分か
非空の兄弟生成語によって分離される．
これは Yoshinaka の多次元文脈の規約に対応する．
特に $\lambda$-自由性から，
任意の非終端記号に由来するすべての終端タプルの成分が
非空であることが帰納的に従う．

以下の標準化では，
Seki らの補題2.2の証明の非消去化段階 (f3) だけを，
Kanazawa が記録したランク保存形式で用いる．
それに成分の並べ替え，空成分の明示的除去，
Yoshinaka の非融合化を組み合わせる
~\cite[proof of Lemma~2.2, item~(f3)]{SekiEtAl1991}
\cite[Section~2.1]{Kanazawa2019Ogden}
\cite[Section~2.3, Lemma~1]{Yoshinaka2011}．
Yoshinaka はさらに，固定した次元と規則ランクについて，
これらの標準変換でサイズ $O(\|G\|)$ の同値な良形文法が得られることを，
第2.3節のExample~2の後で述べている
~\cite[Section~2.3]{Yoshinaka2011}．

ファンアウトと規則ランクは異なる文法資源である．
一般の MCFG/LCFRS 表示について，
任意の固定ファンアウトを保存する普遍的な二分化定理はない．
規則ランクの縮小には，ファンアウトの増加が必要となることがある
~\cite{RambowSatta1999,GomezRodriguezEtAl2009}．
したがって，以下の質的再構成は
二分標準形も既知の規則ランク上界も課さずに定義する．

\paragraph{縮約済み文法．}
MCFG の非終端記号依存グラフは，
左辺 $A$ の規則に子 $B$ が現れるとき，
辺 $A\to B$ をもつ．
このグラフで開始記号から到達できる非終端記号を
\emph{到達可能}と呼び，
生成するタプル言語が非空の記号を
\emph{生産的}と呼ぶ．
すべての非終端記号が到達可能かつ生産的なら，
文法は\emph{縮約済み}である．

\subsection{標準化}

\begin{proposition}[ランクを保存する良形標準化]
\label{prop:good-normalization}
$G_{\mathrm{in}}$ を規則ランクが高々
$r_{\max}$ の有限標準 $f$-MCFG とする．
このとき，規則ランクが高々 $r_{\max}$ の有限良形
$f$-MCFG $G_{\mathrm{good}}$ が存在して，
\[
  L(G_{\mathrm{good}})
  =L(G_{\mathrm{in}})\setminus\{\eps\}
\]
を満たす．
$f,r_{\max}$ を固定すれば，
$\|G_{\mathrm{good}}\|
 =O_{f,r_{\max}}(\|G_{\mathrm{in}}\|)$ である．
$\eps\notin L(G_{\mathrm{in}})$ なら
両文法は同じ言語を生成する．
\end{proposition}

\begin{proof}
まず到達不能または非生産的な非終端記号を除き，
Seki らの補題2.2の証明の非消去化段階 (f3) を，
Kanazawa のランク保存形式で適用する
~\cite[proof of Lemma~2.2, item~(f3)]{SekiEtAl1991}
\cite[Section~2.1]{Kanazawa2019Ogden}．
非終端記号に成分置換を添字として付けることで
文法を非置換にする．
次に，付録~\ref{app:normalization-details} で明示するように，
終端導出によって実現可能な非空成分パターンで非終端記号を装飾し，
空成分を射影で除去する．
これにより開始記号の $\eps$ が取り除かれ，
残るすべての出力テンプレートが $\lambda$-自由になる．
最後に Yoshinaka の非融合化を適用する
~\cite[Section~2.3, Lemma~1]{Yoshinaka2011}．
これらの操作はファンアウトも規則ランクも増加させない．

$f,r_{\max}$ が固定されていれば，
成分パターンの選択による記号の装飾は高々 $2^f$ 通り，
各規則での子のパターン選択は高々 $2^{fr_{\max}}$ 通りであり，
置換による装飾は高々 $f!$ 個のコピーにすぎない．
Yoshinaka の変換も，
これらの固定パラメータに関して線形サイズである．
したがって総符号長は
$O_{f,r_{\max}}(\|G_{\mathrm{in}}\|)$ のままである．
導出の対応と空成分に関する構成の詳細は
付録~\ref{app:normalization-details} に記録する．
\end{proof}

この命題はサイズ・ファンアウト・規則ランクを制御するが，
任意の入力文法の規則厚みと，
標準化後の厚みが一致するとは主張しない．
そのため厚みに依存する結果はまず良形表示について証明する．
消去規則を含む縮約済み $\lambda$-自由入力表示については，
命題~\ref{prop:structural-good-thickness} が，
元文法のサイズと厚みの多項式尺度を
後で用いる良形文法へ転移する．

\subsection{タプル文脈と代入可能性}

\paragraph{文をつくる文脈．}
\emph{アリティ $d$ の文文脈}とは，
\[
 E=v_0\blank_1v_1\cdots v_{d-1}\blank_dv_d
\]
という形の語で，各番号付き穴がちょうど一度ずつ出現し，
$v_0,v_d\in\Sigma^*$ かつ
内部の区切り語 $v_1,\ldots,v_{d-1}$ は
すべて $\Sigma^+$ に属するものをいう．
こうした文脈の集合を $\mathcal E_d$ と書く．
これは Yoshinaka の多次元文脈の規約である
~\cite[Section~3.1]{Yoshinaka2011}．
タプル $\tuple w=(w_1,\ldots,w_d)$ に対し，
$E[\tuple w]$ は各 $\blank_i$ を $w_i$ で埋めて得られる語とし，
\[
 |E|_{\mathrm{fix}}:=\sum_{i=0}^d|v_i|
\]
とおく．
$E[\tuple w]\in L$ であるとき，
文脈 $E$ は言語 $L$ において
$\tuple w$ を\emph{受理する}という．

\begin{definition}[標本語におけるタプルの出現]
\label{def:tuple-occurrence}
$K\subseteq\Sigma^*$，$w=a_1\cdots a_n\in K$ とする．
語 $w$ の切断位置に
$\cutpos{0}a_1\cutpos{1}\cdots a_n\cutpos{n}$
と番号を付ける．
$0\le i<i'\le n$ に対して
$w[i:i']:=a_{i+1}\cdots a_{i'}$ と書く．

語 $w$ における\emph{アリティ $d$ のタプル出現}とは，
左から右に並ぶ非空区間の列
\[
 0\le i_1<i_1'<i_2<i_2'<\cdots<i_d<i_d'\le n
\]
である．
これによりタプル
$\tuple x=(w[i_1:i_1'],\ldots,w[i_d:i_d'])\in(\Sigma^+)^d$
と，各区間を $\blank_1,\ldots,\blank_d$ に置き換えた文脈
$E\in\mathcal E_d$ が決まる．
出現を $(E,\tuple x)$ と同一視する．
例えば $\mathtt{abc}$ の中の
$\mathtt a$ と $\mathtt c$ を選べば，
タプル $(\mathtt a,\mathtt c)$ と
文脈 $\blank_1\mathtt b\blank_2$ が得られる．
\end{definition}

$L\subseteq\Sigma^*$，
$\tuple x\in(\Sigma^*)^d$ に対して，
\emph{$d$-タプル分布}を
\[
 \D_L^{(d)}(\tuple x):=
 \{E\in\mathcal E_d:E[\tuple x]\in L\}
\]
と定義する．アリティが明白なら上付きの $(d)$ を省略する．

\begin{definition}[$(f,h)$-タプル代入可能性]
\label{def:tuple-substitutable}
$f\ge1$ と有限モノイド準同型
$h:\Sigma^*\to M$ を固定する．
言語 $L\subseteq\Sigma^*$ が
\emph{$(f,h)$-タプル代入可能}であるとは，
任意の $1\le d\le f$ と
$\tuple x,\tuple y\in(\Sigma^+)^d$ に対し，
\[
 h^{(d)}(\tuple x)=h^{(d)}(\tuple y),\qquad
 \D_L^{(d)}(\tuple x)\cap\D_L^{(d)}(\tuple y)\ne\emptyset
\]
ならば
$\D_L^{(d)}(\tuple x)=\D_L^{(d)}(\tuple y)$
が成り立つことをいう．
したがって，タプルの成分として空語は比較しない．
標的に $\eps$ が属するときは，開始記号で別途扱う．
これは多次元代入可能性に関する Yoshinaka の
非空成分の規約である
~\cite[Section~3.1 and footnote~3]{Yoshinaka2011}．
CFG の場合は Clark~\cite{Clark2013} も参照されたい．
この制約はファンアウト1ですでに重要である．
空語を比較対象に許せば，
通常の意味では代入可能な $\mathtt a^+$ でさえ，
$\mathtt a$ と $\eps$ が文脈を共有する一方で分布が異なるため，
条件を満たさなくなる．
\end{definition}

ファンアウト1では文文脈がすべて
$u\blank_1v$ という形なので，
$\D_L^{(1)}(x)$ は通常の両側分布
\[
 \D_L(x)=\{(u,v):uxv\in L\}
\]
と自然に同一視できる．
よって定義~\ref{def:tuple-substitutable} は，
外部の規約を仮定することなく，
非空因子に対する通常の固定 $h$ の
両側代入可能性条件に特殊化する．

\paragraph{型付けと分布による同値関係．}
$\tuple x,\tuple y\in(\Sigma^+)^d$ に対し，
\[
 \tuple x\equiv_{L,h}^d\tuple y
 \iff
 h^{(d)}(\tuple x)=h^{(d)}(\tuple y)
 \ \text{かつ}\
 \D_L^{(d)}(\tuple x)=\D_L^{(d)}(\tuple y)
\]
と書く．
$L,h,d$ を固定すると，これは同値関係である．
アリティが明らかな場合は上付き $d$ を省略する．

\begin{definition}[標的言語クラス]
\label{def:class-cmcf}
$f\ge1$ と有限モノイド準同型 $h$ を固定する．
$\Cmcf{f}{h}$ は，
$(f,h)$-タプル代入可能な $f$-MCFL 全体のクラスである．
\end{definition}

$L\in\Cmcf{f}{h}$ において，
$h$ が境界型付けであるとき，
$L$ を\emph{境界型付き}と呼ぶ．

$h\preceq h'$ ならば
$\Cmcf{f}{h}\subseteq\Cmcf{f}{h'}$ である．
より細かい $h'$-型が等しければ
より粗い $h$-型も等しいため，
$h$ の下で使える代入の結論を適用できるからである．
ファンアウト1では，先述の両側固定 $h$ 条件の
文脈自由部分に一致する．

任意の $f\ge1$ について
\[
 \Cmcf{f+1}{h}\cap f\text{-}\mathrm{MCFL}
 \subseteq\Cmcf{f}{h}
\]
が成り立つ．
ファンアウト $f$ の条件で調べるのは
アリティ $1,\ldots,f$ だけだからである．
検査するアリティを増やすと代入可能性が失われる現象は，
自明型付けでも既に現れる．
Yoshinaka は $\mathsf{SL}(p,r)$ と
$\mathsf{SL}(p+1,r)$ が一般に比較不能であることを指摘し，
1次元では代入可能だが2次元では代入可能でない言語
$L_{\mathrm{reverse}}=\{w\#w^R\}$
を与えている
~\cite[Section~4.1 and Example~5]{Yoshinaka2011}．
本稿の有限型付けが追加するのは，
命題~\ref{prop:strong-fanout-separation} のより強い現象である．
すなわち，ある有限型付けの下ではファンアウト1で学習可能な標的が，
\emph{いかなる}有限型付けでも
ファンアウト2の条件を満たさないことがある．
したがって $f$ は表現能力の上界というだけでなく，
学習バイアスの一部である．
学習対象クラスの $f$ に関する単調性は主張しない．

任意の $\mathsf{Acc}\subseteq M$ について，
$h^{-1}(\mathsf{Acc})\in\Cmcf{f}{h}$ である．
実際，この言語は正則であり，
対応する成分の $h$-型が等しければ，
任意の文文脈に入れたときの全体の $h$-値も等しい．
したがってタプル分布が等しくなる．

\begin{proposition}[交差とアーベル群フィルタ]
\label{prop:abelian-coset-filter}
$f$ と $h$ を固定する．
\begin{enumerate}[label=(\roman*),leftmargin=*]
\item $L_1,L_2$ がともに $(f,h)$-タプル代入可能なら，
$L_1\cap L_2$ も $(f,h)$-タプル代入可能である．
\item $L$ が $(f,h)$-タプル代入可能とする．
$\mathbb A$ をアーベル群，
$\varphi:\Sigma^*\to\mathbb A$ をモノイド準同型，
$\mathcal H\le\mathbb A$，
$\gamma\in\mathbb A$ とすると，
$L\cap\varphi^{-1}(\gamma+\mathcal H)$
は $(f,h)$-タプル代入可能である．
\end{enumerate}
得られた交差言語が $f$-MCFL なら，
$\Cmcf{f}{h}$ に属する．
\end{proposition}

\begin{proof}
(i) について，
同じ型の $\tuple x,\tuple y$ が
$E[\tuple x],E[\tuple y]\in L_1\cap L_2$
を満たす共通文脈 $E$ をもつなら，
各 $L_i$ で両者のタプル分布は等しく，
その交差でも等しい．

(ii) については，まず
$\varphi^{-1}(\gamma+\mathcal H)$ が
$(f,h_{\mathrm{triv}})$-タプル代入可能であることを示す．
共通文脈が $\tuple x,\tuple y$ を受理するなら，
\[
 \sum_i(\varphi(y_i)-\varphi(x_i))\in\mathcal H
\]
である．
可換性により，任意の別の文文脈 $E'$ に対して，
差
$\varphi(E'[\tuple y])-\varphi(E'[\tuple x])$
は同じ $\mathcal H$ の元となる．
したがって剰余類への所属は両方向で保存される．
自明型付けは $h$ より粗いので，
この剰余類言語は $(f,h)$-タプル代入可能でもある．
(i) を適用すればよい．
\end{proof}

直ちに次が従う．
$\mathsf{Acc}\subseteq M$，
アーベル群 $\mathbb A$，
準同型 $\varphi:\Sigma^*\to\mathbb A$，
部分群 $\mathcal H\le\mathbb A$，
$\gamma\in\mathbb A$ に対し，
\[
 h^{-1}(\mathsf{Acc})
 \cap\varphi^{-1}(\gamma+\mathcal H)
\]
という形の任意の $f$-MCFL は
$\Cmcf{f}{h}$ に属する．
したがって，一つの型付けの下で
アーベル群の制約を変化させれば，
標的言語族を得られる．
これらの因子化は学習器には与えない．

\section{正例データからの厳密な再構成}
\label{sec:learner}

We define the finite-sample reconstruction operator, prove soundness, and then
prove completeness using a normalized target grammar refined by yield type.
The target grammar is used only to select a finite characteristic sample.
Since every observed tuple component is nonempty, the rank of an enumerated
composition witness is bounded by the length of the sample word containing it.

\subsection{Set-Driven Reconstruction and Sample-Bounded Rank}


Following Definition~\ref{def:tuple-occurrence}, enumerate all tuple occurrences
of arity at most $f$ in the sample $K$.  The reconstruction uses one nonterminal
$[\tuple x]$ for every observed tuple value $\tuple x\in(\Sigma^+)^d$, together
with a fresh start symbol $\widehat S$.  It is not told which occurrences
correspond to target nonterminals or rules; all admissible segmentations are
enumerated.

\begin{definition}[Composition witness]
\label{def:composition-witness}
Let $r\ge1$.  A \emph{composition witness of rank $r$} in $K$ consists of an
observed parent occurrence $(E_{\mathrm{par}},\tuple z)$ of arity $e\le f$ together with a
segmentation of the parent components into terminal gaps and nonempty labelled
intervals.  The labels are variables $\xi_j^{\kappa}$ with
$1\le j\le r$ and $1\le\kappa\le d_j\le f$; every variable occurs exactly once,
and for each fixed $j$ the variables
$\xi_j^1,\ldots,\xi_j^{d_j}$ occur in this order when the parent components are
scanned from left to right.  No output component may contain adjacent variables
$\xi_j^{\kappa}\xi_j^{\kappa+1}$ of the same child.  The labelled intervals
induce child tuples $\tuple u_j=(u_{j,1},\ldots,u_{j,d_j})$ and determine a
unique linear template $\rho$ satisfying
\[
  \tuple z=\rho(\tuple u_1,\ldots,\tuple u_r).
\]
\end{definition}

\begin{lemma}[Materialization of child tuples]
\label{lem:composition-witness-materialization}
Every child tuple induced by a composition witness is itself an observed tuple
occurrence, and the induced template $\rho$ is good.
\end{lemma}

\begin{proof}
Because the parent is an observed tuple occurrence, the components of each
$\tuple u_j$ occur in their original left-to-right order in the same sample
word.  If two consecutive components of the same child lie in one parent
component, the definition of a composition witness excludes their direct
adjacency.  They are therefore separated either by nonempty terminal material
or by one or more nonempty labelled intervals belonging to other children.
If they lie in different parent components, they are separated by the
nonempty internal material of the parent sentence context.  Hence the
components of each $\tuple u_j$ satisfy the separation condition for an
observed tuple occurrence.

Every variable occurs exactly once, the components of each child occur in
order, and adjacent consecutive components of one child are excluded.
Therefore $\rho$ is linear, nondeleting, nonpermuting, and nonmerging.
Finally every parent component is nonempty, so every output component of
$\rho$ is nonempty.  Thus $\rho$ is good.
\end{proof}

\begin{lemma}[Enumeration and sample-bounded rank]
\label{lem:general-witness-enumeration}
Fix $f$ and a finite sample $K$, and put $n_K:=\max\{1,\|K\|_+\}$.  If a
rank-$r$ composition witness has parent tuple $\tuple z$ occurring in a sample
word $w$, then
\[
  r\le\|\tuple z\|_1\le |w|\le\|K\|_+.
\]
All composition witnesses are effectively enumerable.  Under a rule-rank cap
$r_{\max}$, those of rank at most $r_{\max}$ are enumerable in
$n_K^{O(f r_{\max})}$ time and space.
\end{lemma}

\begin{proof}
Every child tuple has at least one nonempty component.  Linearity and
nondeletion make the child components contribute disjointly to the parent, so
$\|\tuple z\|_1\ge\sum_j\|\tuple u_j\|_1\ge r$; the remaining inequalities
follow because the parent components are disjoint intervals of $w$.

An arity-$d$ tuple occurrence is determined by $2d$ cuts.  For a fixed rank
$r$, a candidate segmentation has at most $fr$ labelled child intervals, whose endpoints are chosen among the $|w|+1\le n_K+1$ cut positions of
the parent sample word.  For $r\le r_{\max}$, assignments to parent components and
terminal gaps contribute only a factor depending on $f$ and $r_{\max}$, and
the nonpermuting and nonmerging conditions are directly checkable.  Summing
over $r\le r_{\max}$ gives the stated bound.  Without a cap, the displayed
rank inequality restricts the effective enumeration to
$1\le r\le\|K\|_+$.
\end{proof}

\begin{definition}[Capped and uncapped reconstruction operators]
\label{def:reconstruction}
Fix $f$ and $h$.  For $r_{\max}\ge1$, let
$\mathcal B_{f,h}^{\le r_{\max}}(K)$ be the grammar
$\widehat G_{\le r_{\max}}(K)$ with a fresh start symbol $\widehat S$, one nonterminal $[\tuple x]$ for each nonempty tuple value $\tuple x$ observed in $K$,
and the following rules:
\begin{enumerate}[leftmargin=*]
\item[\textnormal{(Start)}] for every $w\in K\cap\Sigma^+$,
      $\widehat S\to[(w)]$, and if $\eps\in K$, also
      $\widehat S\to\eps$;
\item[\textnormal{(Const)}] for every observed tuple, the constant rule
      $[\tuple u]\to\tuple u$;
\item[\textnormal{(Comp)}] for every composition witness of rank
      $1\le r\le r_{\max}$ with parent tuple $\tuple z$, template $\rho$,
      and induced child tuples $\tuple u_1,\ldots,\tuple u_r$, the rule
      \[
        [\tuple z]\to\rho([\tuple u_1],\ldots,[\tuple u_r]);
      \]
\item[\textnormal{(Link)}] for observed nonempty tuples
      $\tuple x,\tuple y$ of the same arity $d$, the identity-template rule
      $[\tuple x]\to[\tuple y]$ whenever
      $h^{(d)}(\tuple x)=h^{(d)}(\tuple y)$ and some $E\in\mathcal E_d$
      satisfies $E[\tuple x],E[\tuple y]\in K$.
\end{enumerate}
If $K\ne\emptyset$, define
$\mathcal B_{f,h}(K):=\widehat G(K):=
\widehat G_{\le\|K\|_+}(K)$; for $K=\emptyset$, let
$\mathcal B_{f,h}(\emptyset)$ be the grammar with no productive start rule.
Fixed encoding orders make both reconstruction operators deterministic and
set-driven.
\end{definition}

\paragraph{Exact relation to Yoshinaka's finite-sample grammar.}
The finite typing modifies Yoshinaka's constructor only at substitution links.
\begin{proposition}[Typed guard on Yoshinaka's reconstruction]
\label{prop:yoshinaka-type-guard}
Fix a finite positive sample $K$, fan-out bound $f$, and rank cap
$r_{\max}\ge1$.  Under the trivial typing,
$\mathcal B_{f,h_{\mathrm{triv}}}^{\le r_{\max}}(K)$ agrees, up to notation,
with Yoshinaka's hypothesis grammar $G(K)$ for $\mathcal A(f,r_{\max})$
\cite[Section~4.1]{Yoshinaka2011}, except that the present construction handles
$\eps\in K$ by a separate direct start rule.  For an arbitrary finite typing
$h$, $\mathcal B_{f,h}^{\le r_{\max}}(K)$ is obtained from that grammar by
deleting exactly those Type~II substitution rules whose source and target
tuples have different componentwise $h$-types.  The same statement holds for
the uncapped reconstruction and Yoshinaka's $\mathcal A(f,*)$.
\end{proposition}

\begin{proof}
The observed tuple values are the same nonterminal indices in both
constructions.  Rule~(Start) is Yoshinaka's Type~III rule, and Rule~(Link) is
Type~II with the additional componentwise type test.

It remains to compare Type~I with Rules~(Const) and~(Comp).  A rank-zero
Type~I rule is exactly Rule~(Const).  For positive rank, suppose Yoshinaka has
\[
 [\tuple z]\to\rho([\tuple u_1],\ldots,[\tuple u_r])
 \qquad\text{with}\qquad
 \tuple z=\rho(\tuple u_1,\ldots,\tuple u_r),
\]
where all displayed tuple values are observed and $\rho$ is good.  Choose an
observed occurrence of $\tuple z$ in a sample word.  The displayed equality
locates every component of every $\tuple u_j$ as a labelled interval inside
that occurrence.  Linearity and nondeletion use each interval once,
nonpermutation puts the components of each child in order, and nonmerging
prevents direct adjacency of consecutive components of one child.  Hence these
intervals form a composition witness, so Rule~(Comp) contains the same rule.
Conversely every composition witness is a good algebraic decomposition among
observed tuple values and therefore is a Type~I rule.  This also matches
Yoshinaka's counting proof, where each Type~I rule is enumerated as a
factorization of a single sample word
\cite[Lemma~6]{Yoshinaka2011}.

Thus trivial typing gives the same Type~I--III rule set, apart from the direct
$\eps$ start convention.  For general $h$, Rules~(Start), (Const), and (Comp)
are unchanged, and only Type~II/Rule~(Link) is filtered by componentwise type.
Finally, in the uncapped construction every composition witness has rank at
most the length of the sample word containing it, so allowing the cap
$\|K\|_+$ introduces no Type~I rule beyond the sample-length cap used in
Yoshinaka's $\mathcal A(f,*)$.
\end{proof}

\begin{example}[Why the equal-type premise matters]
\label{ex:type-guard-purpose}
Let $L=\{\mathtt a^n\mathtt b^n\mathtt c^n:n\ge1\}$ and
$K=\{\mathtt{abc},\mathtt{aabbcc}\}\subseteq L$.  The observed factors
$(\mathtt b)$ and $(\mathtt{abbc})$ share the sample context
$\mathtt a\blank_1\mathtt c$.  With trivial typing, Rule~(Link) therefore
connects their hypothesis nonterminals.  The composition witness
$[\mathtt{aabbcc}]\to\mathtt{aa}[\mathtt b]\mathtt{bcc}$ can then use that
link to derive $\mathtt{aaabbcbcc}\notin L$.  Under $h_\Sigma^{1,1}$ the
link is blocked because $h_\Sigma^{1,1}(\mathtt b)=\mathtt b$ while
$h_\Sigma^{1,1}(\mathtt{abbc})=\langle\mathtt a,\mathtt c\rangle$.
This illustrates the local role of the equal-type premise.
\end{example}

The reconstruction is sample-consistent:
\[
 K\subseteq L\bigl(\mathcal B_{f,h}^{\le r_{\max}}(K)\bigr),
 \qquad
 K\subseteq L\bigl(\mathcal B_{f,h}(K)\bigr).
\]
Indeed, each nonempty $w\in K$ has the direct derivation
$\widehat S\to[(w)]\to(w)$ by Rules~(Start)--(Const), while
$\eps\in K$ gives the direct start rule $\widehat S\to\eps$.

These $\mathcal B$ operators are set-driven constructors, not sequential TxtEx
learners: recomputation on growing samples need not stabilize syntactically
(e.g. for successively longer $\mathtt a^n\mathtt c\mathtt b^n$).  The conservative learner below
rebuilds only when a new positive datum is rejected.

Both reconstruction operators produce finite standard MCFGs of fan-out at
most $f$ and are effectively constructible from $K,f,h$.  There are only
finitely many observed nonempty tuple occurrences.  Rule~(Link) is obtained by
grouping occurrences by concrete context and componentwise type, while
Lemma~\ref{lem:general-witness-enumeration} enumerates Rule~(Comp).  By
Lemma~\ref{lem:general-witness-enumeration}, every rule of $\mathcal B_{f,h}(K)$
has rank at most $\|K\|_+$ when $K\ne\emptyset$.

The concrete separation examples in Section~\ref{sec:comparison} also serve as
small instances of the tuple-nonterminal and composition mechanisms used here.

\subsection{Soundness}
\label{sec:soundness}

\begin{lemma}[Induced child context]
\label{lem:filling-identity}
Let $\rho$ be the template of a good rule, let
$E\in\mathcal E_e$ be a sentence context for its output tuple, and fix tuples of the arities required by the child nonterminals, with nonempty
components, for all children except child $j$.
Replacing every variable of child $j$ by its named hole and every other child
variable by the corresponding fixed component produces a context
$E_j\in\mathcal E_{d_j}$
such that
\[
 E_j[\tuple u_j]=E[\rho(\tuple u_1,\ldots,\tuple u_r)].
\]
If the sibling tuples are changed, the same construction gives the
corresponding induced context $E_j$ for that new sibling choice; the context
itself need not be identical.
\end{lemma}

\begin{proof}
Linearity and nondeletion make every variable of child $j$ occur exactly once,
and nonpermutingness puts $\xi_j^1,\ldots,\xi_j^{d_j}$ in that order.  Between two
consecutive variables of this child, nonmerging ensures either a terminal
symbol, a variable of another child, or a boundary between two output
components.  In the first case the separator is visibly nonempty; in the
second it becomes a nonempty sibling component; in the third the corresponding
internal separator of $E$ is nonempty.  Replacing the distinguished variables
by $\blank_1,\ldots,\blank_{d_j}$ therefore gives a context in
$\mathcal E_{d_j}$, and filling the holes is exactly composition with $\rho$
and then with $E$.
\end{proof}

\begin{lemma}[Composition preserves equivalence]
\label{lem:witnessed-composition}
Let $L\subseteq\Sigma^*$, let $\rho$ be the template of a good rank-$r$ rule, and let
$\tuple x_j,\tuple u_j$ have the arities required by its children and have
nonempty components.  If
$\tuple u_j\equiv_{L,h}\tuple x_j$ for every $j$, then
\[
 \rho(\tuple u_1,\ldots,\tuple u_r)
 \equiv_{L,h}
 \rho(\tuple x_1,\ldots,\tuple x_r).
\]
\end{lemma}

\begin{proof}
Homomorphic template evaluation gives equality of the two parent $h$-types,
and goodness of $\rho$ makes every parent component nonempty.  It remains to
compare distributions.  Let $E$ be an arbitrary sentence context for the
parent tuple.  If
$E[\rho(\tuple x_1,\ldots,\tuple x_r)]\in L$, replace the children one at
a time.  At step $j$, Lemma~\ref{lem:filling-identity} supplies the induced
sentence context for child $j$ under the current sibling choice.  Equality of
the child distributions transfers acceptance from $\tuple x_j$ to
$\tuple u_j$.  After all replacements,
$E[\rho(\tuple u_1,\ldots,\tuple u_r)]\in L$.  Reversing the replacements
gives the converse implication.  Since $E$ was arbitrary, the two parent
tuples have equal distributions.
\end{proof}

\begin{proposition}[Soundness of every capped hypothesis]
\label{prop:reconstruction-soundness}
Let $L$ be $(f,h)$-tuple-substitutable and $K\subseteq L$.  For every
$r_{\max}\ge1$, if $[\tuple x]\Rightarrow^*\tuple u$ in
$\widehat G_{\le r_{\max}}(K)$, then every component of $\tuple u$ is nonempty
and $\tuple u\equiv_{L,h}\tuple x$.  Consequently
\[
 L\bigl(\mathcal B_{f,h}^{\le r_{\max}}(K)\bigr)\subseteq L,
 \qquad L\bigl(\mathcal B_{f,h}(K)\bigr)\subseteq L.
\]
\end{proposition}

\begin{proof}
Induct on derivation height.  Constant rules are immediate.  For a Rule~(Link) instance $[\tuple x]\to[\tuple y]$, its defining common sample context and
equal $h$-type give $\tuple x\equiv_{L,h}\tuple y$; transitivity with the
induction hypothesis proves the claim.  For a composition-witness rule, the
induction hypotheses give nonempty child tuples equivalent to the witnessed
children.  Every parent component of the witnessed tuple was nonempty; under
the same nondeleting template it remains nonempty after replacing child
components by other nonempty strings.  Lemma~\ref{lem:witnessed-composition}
then gives the parent equivalence.

A nonempty start derivation begins with $\widehat S\to[(w_0)]$ for some
$w_0\in K\cap\Sigma^+$, so the empty unary context places the derived word in
$L$.  The direct rule $\widehat S\to\eps$ is present only when
$\eps\in K\subseteq L$.  The $K=\emptyset$ case is trivial, and the general
hypothesis is the capped case with $r_{\max}=\|K\|_+$ whenever $K\ne\emptyset$.
\end{proof}


\subsection{Yield-Typed Target Refinement and Finite Witnesses}
\label{sec:refinement}

\begin{definition}[Template yield type]
\label{def:output-map}
Let
\(R:A\to\rho(B_1,\ldots,B_r)\)
be a nondeleting rule with output template
\(\rho=(\alpha_1,\ldots,\alpha_{d_0})\), and put \(d_0=\fan(A)\) and
\(d_j:=\fan(B_j)\), and choose
\(\mathbf t_j\in M^{d_j}\) for each \(1\le j\le r\).

On the template alphabet
\(
  \Gamma_\rho
  :=\Sigma\uplus\biguplus_{j=1}^r
  \{\xi_j^1,\ldots,\xi_j^{d_j}\}
\),
define
\(\evalmap_{\mathbf t_1,\ldots,\mathbf t_r}\colon\Gamma_\rho\to M\)
by
\(\evalmap_{\mathbf t_1,\ldots,\mathbf t_r}(a):=h(a)\)
for \(a\in\Sigma\), and
\(\evalmap_{\mathbf t_1,\ldots,\mathbf t_r}(\xi_j^{\kappa}):=(\mathbf t_j)_\kappa\).

Let
\(
  \widehat{\evalmap}_{\mathbf t_1,\ldots,\mathbf t_r}
  \colon\Gamma_\rho^*\to M
\)
be its unique extension to a monoid homomorphism.
Define
\[
  \out_\rho(\mathbf t_1,\ldots,\mathbf t_r)
  :=
  \bigl(
    \widehat{\evalmap}_{\mathbf t_1,\ldots,\mathbf t_r}(\alpha_1),
    \ldots,
    \widehat{\evalmap}_{\mathbf t_1,\ldots,\mathbf t_r}(\alpha_{d_0})
  \bigr)
\in M^{d_0}.
\]
For rank zero, \(\out_\rho\) is simply
\((h(\alpha_1),\ldots,h(\alpha_{d_0}))\), the componentwise \(h\)-type of the
constant template.
\end{definition}

The refinement below is ordinary finite nonterminal splitting: compositionality
$h(uv)=h(u)h(v)$ makes each typed target nonterminal uniform in componentwise
$h$-type, which is exactly the guard needed for substitutability.


\begin{definition}[Yield-type refinement]
\label{def:output-refinement}
Given a finite nondeleting MCFG \(G=(V,\Sigma,P,S,\fan)\) and a finite
typing \(h\), the \emph{full yield-type refinement} \(G^h\) has a fresh
start symbol \(\widetilde S\) and typed nonterminals
\(A_{\mathbf t}\) for every \(A\in V\) and
\(\mathbf t\in M^{\fan(A)}\).
Set
\(\fan(\widetilde S):=1\) and
\(\fan(A_{\mathbf t}):=\fan(A)\).

For each \(\eta\in M\), add the start rule
\(\widetilde S\to S_{(\eta)}\).
For each rank-zero rule \(A\to\tuple c\), add
\(A_{h^{(\fan(A))}(\tuple c)}\to\tuple c\).

For every positive-rank rule
\[
R:A\to\rho(B_1,\ldots,B_r),
\qquad
\rho=(\alpha_1,\ldots,\alpha_{d_0}),
\]
and every choice $\mathbf t_j\in M^{\fan(B_j)}$ ($1\le j\le r$), add
\[
  A_{\out_\rho(\mathbf t_1,\ldots,\mathbf t_r)}
  \to
  \rho(B_{1,\mathbf t_1},\ldots,B_{r,\mathbf t_r}).
\]
Here \(B_{j,\mathbf t_j}\) denotes the typed copy of \(B_j\) indexed by
\(\mathbf t_j\).

Finally, the \emph{trimmed refinement} \(\widetilde G_0\) is obtained by
keeping only the typed nonterminals and typed rules that actually occur in some
successful derivation from \(\widetilde S\); the subscript $0$ denotes this
productive/reachable core.  From this point on, a ``typed nonterminal'' or
``typed rule'' means one of \(\widetilde G_0\) unless the full refinement
\(G^h\) is mentioned explicitly.
\end{definition}

\begin{proposition}[Yield-type invariants]
\label{prop:output-invariants}
For \(G^h\) and \(\widetilde G_0\):
\begin{enumerate}[label=(\roman*),leftmargin=*]
\item if \(A_{\mathbf t}\Rightarrow^*\tuple u\), then
\(h^{(\fan(A))}(\tuple u)=\mathbf t\), and erasing the type indices gives a
valid \(G\)-derivation of \(\tuple u\);
\item every terminal derivation
\(A\Rightarrow_G^*\tuple u\)
lifts uniquely to a derivation
\(A_{h^{(\fan(A))}(\tuple u)}\Rightarrow_{G^h}^*\tuple u.\)
In particular, every successful derivation \(S\Rightarrow_G^* w\) lifts
uniquely through the start rule
\(
  \widetilde S\to S_{(h(w))}
\);
\item \(L(\widetilde G_0)=L(G^h)=L(G)\).
\end{enumerate}
\end{proposition}

\begin{proof}
At a rule application, homomorphic evaluation of the rule template gives
\[
  h^{(\fan(A))}
  \bigl(\rho(\tuple u_1,\ldots,\tuple u_r)\bigr)
  =
  \out_\rho
  \bigl(h^{(\fan(B_1))}(\tuple u_1),\ldots,
        h^{(\fan(B_r))}(\tuple u_r)\bigr).
\]
The rank-zero case is immediate.  Induction on derivation height proves~(i),
and the same identity gives the unique terminal-derivation lifting in~(ii).  Erasing the
type indices gives one language inclusion.  Conversely, every successful
terminal derivation of \(G\) has the unique typed lifting described in~(ii),
including the uniquely determined start link
\(\widetilde S\to S_{(h(w))}\).  Finally, every successful typed derivation
consists entirely of rules retained by the trimming, so trimming preserves the
generated language.
\end{proof}

Yield typing changes only nonterminal labels, so every typed rule has the
same rule template as its untyped source.  Hence $\lambda$-freeness,
nondeletion, nonpermutation, and nonmerging are all inherited by the trimmed
refinement.

For a typed nonterminal $X=A_{\mathbf t}$, put
$L_X:=\{\tuple u:X\Rightarrow_{\widetilde G_0}^*\tuple u\}$.  Trimming makes
$L_X$ nonempty; if the target grammar is $\lambda$-free, every component
of every $\tuple u\in L_X$ is nonempty.

\begin{lemma}[Canonical contexts for typed nonterminals]
\label{lem:identity-canonical-contexts}
Assume the target grammar is good.  Every typed nonterminal $X$ of arity $d$ has a sentence context
$\chi(X)\in\mathcal E_d$ such that
$\chi(X)[\tuple u]\in L(G)$ for every $\tuple u\in L_X$.  If
$\widetilde S\to X$ belongs to $\widetilde G_0$, one may take $\chi(X)=\blank_1$.
\end{lemma}

\begin{proof}
Choose a successful typed derivation containing an occurrence of $X$.  At the
root, the one-hole context $\blank_1$ is valid.  Follow the unique path from
the root occurrence down to the chosen occurrence of $X$.  At each rule on
this path, fix the productive sibling subderivations used in the successful
derivation and apply Lemma~\ref{lem:filling-identity}.  Goodness of the rule
and nonemptiness of the sibling components ensure that the induced child
context again has nonempty internal separators.  Iterating down the path gives
$\chi(X)\in\mathcal E_d$.  The surrounding typed rule applications depend only on the typed symbol
$X$, so replacing its subtree by any other $X$-derivation remains a valid
derivation in $\widetilde G_0$.  If $\widetilde S\to X$ is a rule of
$\widetilde G_0$, the path has length zero and $\chi(X)=\blank_1$.
\end{proof}

When $L=L(G)\in\Cmcf{f}{h}$, the refinement is also distributionally homogeneous for the sentence-context
family used here.  If $X=A_{\mathbf t}$ has arity $d$ and
$\tuple u,\tuple v\in L_X$, Proposition~\ref{prop:output-invariants} gives
$h^{(d)}(\tuple u)=h^{(d)}(\tuple v)=\mathbf t$, while
Lemma~\ref{lem:identity-canonical-contexts} supplies one sentence context
accepting both.  Hence $(f,h)$-tuple substitutability gives
$\D_L^{(d)}(\tuple u)=\D_L^{(d)}(\tuple v)$.  Thus each typed nonterminal lies
inside one distributional class for $\mathcal E_d$.  This parallels, but is weaker than, the full-multicontext homogeneity used
for congruential MCFGs~\cite{ClarkYoshinaka2016}: our internal separators are
required to be nonempty.  For genuine tuple arity $d\ge2$, the present
sentence-context family still requires nonempty internal separators, so we do
not assert an inclusion into the full multicontext congruential MCFG class;
whether the empty-separator cases can be recovered uniformly remains open.

\begin{definition}[Canonical tuples and contexts]
\label{def:target-canonical}
For every non-start typed nonterminal $X$, fix a \emph{canonical tuple}
$\omega(X)\in L_X$ and a \emph{canonical context} $\chi(X)$ supplied by
Lemma~\ref{lem:identity-canonical-contexts}.  Whenever
$\widetilde S\to X$ belongs to $\widetilde G_0$, require
$\chi(X)=\blank_1$; otherwise choose any canonical context from the lemma.  For
a $\lambda$-free target grammar every component of $\omega(X)$ is nonempty,
so $\|\omega(X)\|_1\ge\fan(X)\ge1$.  Qualitatively a canonical tuple may be the first
tuple in an effective enumeration of finite terminal derivations; the
quantitative sections choose short canonical tuples.  These canonical objects are target witnesses, not inputs to the reconstruction
operator.
\end{definition}

\begin{definition}[Finite target witness set]
\label{def:witness-set}
Let $\Wit(\widetilde G_0)$ consist of the following positive words:
\[
\begin{aligned}
&\chi(X)[\omega(X)]
&&\text{for every non-start typed nonterminal }X,\\
&\chi(X)[\rho(\omega(Y_1),\ldots,\omega(Y_r))]
&&\text{for every non-start positive-rank typed rule }
  X\to\rho(Y_1,\ldots,Y_r),\\
&\chi(X)[\tuple c]
&&\text{for every non-start rank-zero typed rule }X\to\tuple c.
\end{aligned}
\]
\end{definition}

Because $\widetilde G_0$ is finite and all of its symbols and rules occur in
successful derivations, $\Wit(\widetilde G_0)$ is a finite subset of $L(G)$;
when $L(G)\ne\emptyset$ it is nonempty.

\begin{lemma}[Target rule words occur in the enumeration]
\label{lem:exposed-rule-enumerated}
Let $R:X\to\rho(Y_1,\ldots,Y_r)$ be a non-start positive-rank typed rule.
Put $\tuple z_R:=\rho(\omega(Y_1),\ldots,\omega(Y_r))$.  If
$\Wit(\widetilde G_0)\subseteq K$, the witness enumeration contains
\[
  [\tuple z_R]\to\rho([\omega(Y_1)],\ldots,[\omega(Y_r)]).
\]
\end{lemma}

\begin{proof}
The word $\chi(X)[\tuple z_R]\in K$ supplies the parent occurrence.
The segmentation of this sample occurrence induced by the variables of
$\rho$ places every canonical tuple of each child as a tuple occurrence in the same sample word.  $\lambda$-freeness makes all child components nonempty, nonpermutingness
puts each child's components in index order, and nonmerging guarantees a
nonempty separator between consecutive components of each child in the
sample word.  Thus every child canonical tuple is an observed tuple occurrence and the
segmentation is a composition witness; exhaustive enumeration emits the
displayed rule.
\end{proof}


\begin{proposition}[Cardinality of the finite target witness set]
\label{prop:witness-cardinality}
Let $G=(V,\Sigma,P,S,\fan)$ be the normalized target presentation from which
$\widetilde G_0$ is constructed, with fan-out at most
$f$ and rule rank at most $r_{\max}$.  For a fixed finite typing
$h:\Sigma^*\to M$, the target witness set constructed
above may be chosen so that
\[
  |\Wit(\widetilde G_0)|
  \le |V||M|^f+|P||M|^{f r_{\max}}.
\]
In particular, for fixed $f$, $r_{\max}$, and $h$, the number of selected positive
words is linear in $|V|+|P|$.
\end{proposition}

\begin{proof}
The yield-type refinement has at most $|V||M|^f$ non-start typed nonterminals.
The source rules give at most $|P||M|^{f r_{\max}}$ non-start typed rule
instances.  The construction selects at most one word for each such typed
nonterminal and one for each such typed rule instance; fresh start links are
handled separately and require no additional witness word.
\end{proof}

Proposition~\ref{prop:witness-cardinality} controls only the number of selected
words, not their total positive length $\|\Wit(\widetilde G_0)\|_+$.
Polynomial total-length bounds require the additional rank-one or
boundary-determined arguments below.


\subsection{Completeness and Exact Reconstruction}

\label{sec:completeness}

\begin{proposition}[Completeness from a target presentation]
\label{prop:completeness}
This is the finite-typed counterpart of Yoshinaka's rule-witness
completeness argument~\cite[Lemma~5]{Yoshinaka2011}.  Let $G$ be a finite good MCFG of
fan-out at most $f$, let
$L_G:=L(G)$, and construct its trimmed yield-type refinement
$\widetilde G_0$ and witness set above.  Let $r_G$ be the maximum rank of a
non-start typed rule, with $r_G:=0$ if there is none.  For every
finite sample $K$ with $\Wit(\widetilde G_0)\subseteq K$ and every
$r_{\max}\ge\max\{1,r_G\}$,
\[
  L_G\subseteq L\bigl(\mathcal B_{f,h}^{\le r_{\max}}(K)\bigr),
\]
and also $L_G\subseteq L\bigl(\mathcal B_{f,h}(K)\bigr)$.
\end{proposition}

\begin{proof}
For each non-start typed nonterminal $X$, put $\iota(X):=[\omega(X)]$.  If
$\widetilde S\to X$ belongs to $\widetilde G_0$, then $\chi(X)=\blank_1$, so the reconstruction contains
the corresponding start rule.

For a non-start rank-zero typed rule $R:X\to\tuple c$, the two corresponding witness words share $\chi(X)$, and Proposition~\ref{prop:output-invariants} gives the same
componentwise $h$-type to $\omega(X)$ and $\tuple c$.  Hence the reconstruction contains
$\iota(X)\to[\tuple c]$, followed by the constant rule.

For a non-start positive-rank typed rule
$R:X\to\rho(Y_1,\ldots,Y_r)$, put
$\tuple z_R:=\rho(\omega(Y_1),\ldots,\omega(Y_r))$.  Yield typing gives
$h^{(\fan(X))}(\tuple z_R)=h^{(\fan(X))}(\omega(X))$, and the two corresponding witness words share $\chi(X)$, so Rule~(Link)
$\iota(X)\to[\tuple z_R]$ is present.  Lemma~\ref{lem:exposed-rule-enumerated}
provides
\[
 [\tuple z_R]\to\rho(\iota(Y_1),\ldots,\iota(Y_r)).
\]
Induction on a $\widetilde G_0$-derivation therefore simulates every non-start
typed rule whenever the cap is at least $r_G$; the fresh start link was handled
above.  Hence every word of $L_G$ is generated.
For the uncapped reconstruction, each required rule has a witness word in $K$;
Lemma~\ref{lem:general-witness-enumeration} places its rank below $\|K\|_+$.  This
simulation uses neither tuple substitutability nor $K\subseteq L$; those enter
only through soundness.
\end{proof}

\begin{theorem}[Exact reconstruction for a fixed typing]
\label{thm:exact-reconstruction}
Fix $f\ge1$ and a finite homomorphism $h:\Sigma^*\to M$.  For every
$L\in\Cmcf{f}{h}$ there is a finite characteristic sample
$K_L^{\star}\subseteq L$ such that every finite
$K_L^{\star}\subseteq K\subseteq L$ satisfies
\[
  L(\mathcal B_{f,h}(K))=L.
\]
\end{theorem}

\begin{proof}
Put $L^+:=L\setminus\{\eps\}$.  If $L^+=\emptyset$, take $K_L^\star=L$.
Otherwise choose a finite $f$-MCFG presentation of $L$ and apply
Proposition~\ref{prop:good-normalization} to obtain a good presentation of
$L^+$.  Let $\widetilde G_0$ be its trimmed yield-type refinement and set
\[
 K_L^\star:=\Wit(\widetilde G_0)\cup(\{\eps\}\cap L).
\]
The observation following Definition~\ref{def:witness-set} makes this a finite
subset of $L$.  For $K_L^\star\subseteq K\subseteq L$,
Proposition~\ref{prop:completeness} gives every nonempty target word in the
hypothesis, while Proposition~\ref{prop:reconstruction-soundness} gives the
reverse inclusion.  If $\eps\in L$, the direct start rule generates it; if
$\eps\notin L$, soundness forbids it.  The cases $L=\emptyset$ and
$L=\{\eps\}$ are immediate from the reconstruction definition.
\end{proof}

The set-driven map that sends a finite text prefix to
$\mathcal B_{f,h}$ of its positive content is eventually semantically correct:
once the content contains $K_L^\star$, every later reconstruction generates
$L$ by Theorem~\ref{thm:exact-reconstruction}.  Its grammar presentation may
still change as new positive examples enlarge the observed tuple inventory,
which is why TxtEx stabilization is handled by the conservative wrapper below.

\begin{definition}[Conservative sequential learner]
\label{def:conservative-learner}
We use the same conservative update policy as Yoshinaka's
Algorithm~1~\cite[Section~4.1]{Yoshinaka2011}.  Let $H_0$ be the fixed
empty-language hypothesis from the learning model.  The
learner $\mathcal A_{f,h}$ starts from $H_0$.  On a pause symbol it keeps its
current grammar.  On a positive datum $w$, it keeps the current grammar if that
grammar already generates $w$; otherwise, with $K$ the positive content seen
so far, it replaces the hypothesis by $\mathcal B_{f,h}(K)$.
\end{definition}

\begin{corollary}[Gold/TxtEx identification]
\label{cor:gold-explanatory}
For every finite-monoid homomorphism $h:\Sigma^*\to M$ and every $f\ge1$,
the conservative learner $\mathcal A_{f,h}$ TxtEx-identifies
$\Cmcf{f}{h}$ from positive data.
\end{corollary}

\begin{proof}
MCFG membership is decidable~\cite{SekiEtAl1991}, so the learner is
computable.  Sample consistency gives the invariant that every current
hypothesis contains every positive example seen so far: a rebuild contains the
whole current content, and a stage without a rebuild leaves a hypothesis that
already contains the new datum.  Proposition~\ref{prop:reconstruction-soundness}
makes every rebuilt hypothesis sound.

Fix a stage after the characteristic sample from
Theorem~\ref{thm:exact-reconstruction} has appeared.  If the current sound
hypothesis is already exact, every later positive datum is generated and no
further rebuild occurs.  Otherwise choose $w\in L$ missing from the current
hypothesis.  The text eventually presents $w$, which forces a rebuild from
positive content still containing the characteristic sample; that rebuild is
exact by Theorem~\ref{thm:exact-reconstruction}.  From then on no further
change is possible.  The learner is generally presentation-order-sensitive and
need not be set-driven.
\end{proof}

\subsection{Polynomial-Time Updates at Bounded Rule Rank}
\label{sec:poly-time}

For fixed $f,r_{\max},h$, the construction is polynomial
in the observed sample, as in Yoshinaka's
bounded-rank learner~\cite[Proposition~4
and Section~4.3]{Yoshinaka2011}.
We encode observed tuples by cut positions and
sample-word indices, and templates by ordered
variables and terminal factors.

\begin{theorem}[Polynomial construction under a rule-rank cap]
\label{thm:poly}
Fix $f,r_{\max}\ge1$ and a finite homomorphism $h$.
From a finite sample $K$ the complete encoding of
$\widehat G_{\le r_{\max}}(K)$ is computable in
time $(1+\|K\|_+)^{O(fr_{\max})}$.
Its fan-out is at most $f$ and rule rank at
most $r_{\max}$.
\end{theorem}

\begin{proof}
With $n_K=\max\{1,\|K\|_+\}$, cut positions
enumerate all tuple occurrences in
$n_K^{O(f)}$ time; fixed-monoid substring
products give their types. Grouping by
context and type enumerates Rule~(Link).
Lemma~\ref{lem:general-witness-enumeration}
enumerates Rule~(Comp) and its encodings
in $n_K^{O(fr_{\max})}$ time.
The remaining rules have rank at most one,
and deduplication leaves a polynomial-size
complete grammar encoding.
\end{proof}

\begin{corollary}[Polynomial-time conservative learning at bounded rank]
\label{cor:capped-gold-poly}
For fixed $f,r_{\max},h$, the policy in
Definition~\ref{def:conservative-learner}
using $\mathcal B^{\le r_{\max}}_{f,h}$
TxtEx-identifies
$\Cmcf{f}{h}\cap\mathsf{MCFL}(f,r_{\max})$
in time polynomial in the accumulated
positive-sample size per processed datum.
\end{corollary}

\begin{proof}
Rank-preserving normalization
(Proposition~\ref{prop:good-normalization})
and soundness and completeness
(Propositions~\ref{prop:reconstruction-soundness}
and~\ref{prop:completeness}) provide the
characteristic sample needed in the proof
of Corollary~\ref{cor:gold-explanatory}.
Theorem~\ref{thm:poly} bounds rebuilds.
Uniform MCFG membership for bounded
fan-out and rank is polynomial jointly
in the grammar encoding and input length
\cite[Proposition~4]{Yoshinaka2011};
the parser saturates finite unary closures.
\end{proof}

At fan-out one every CFG admits a rank-two
presentation, with $\eps$ handled separately.
Thus $\mathcal A^{\le2}_{1,h}$ covers
$\Cmcf{1}{h}$ without implying polynomial
characteristic-data length at rank two.

\paragraph{Relation to Yoshinaka's general framework.}
At bounded rule rank the present typed class is exactly
the substitutable class of a finite-sort signature,
and our reconstruction generates the same language as
Yoshinaka's substructure-indexed $\mathrm{SUBST}_{\mathrm P}$
on each sample. These qualitative conclusions are
instances of the general framework
\cite{Yoshinaka2015General}; precise statements,
context-function correspondence, and proofs appear
in Appendix~\ref{sec:general-framework}.

\section{Characteristic-Data Bounds}
\label{sec:char-data}

Under a rule-rank cap, Proposition~\ref{prop:witness-cardinality} already
gives polynomially many characteristic words; this section controls their
total length.

\begin{center}
\small
\begin{tabularx}{\textwidth}{>{\raggedright\arraybackslash}p{0.27\textwidth}
                              >{\raggedright\arraybackslash}p{0.35\textwidth}
                              X}
\hline
Condition & Characteristic-data control & Source \\
\hline
arbitrary $h$, bounded rank
  & polynomial in $\|G\|$ and refined thickness $\theta_{\widetilde G_0}$
  & Lemma~\ref{lem:typed-thickness-baseline} \\
rule rank one
  & polynomial in $\|G\|$ alone
  & Theorem~\ref{thm:rankone-polydata} \\
$h\preceq h_\Sigma^{k,\ell}$, bounded rank
  & $O(\|G\|^2\bar\theta_G)$
  & Theorem~\ref{thm:boundary-polydata} \\
window plus finite abelian group
  & $O(\|G\|^3\bar\theta_G)$
  & Theorem~\ref{thm:window-abelian-polydata} \\
\hline
\end{tabularx}
\end{center}

The first row is the unconditional baseline for general $h$.
Full yield-type refinement can make its thickness exponentially larger than
the original presentation parameters
(Proposition~\ref{prop:typed-thickness-exponential}), so that route alone
gives no original-thickness bound.  Rank one removes the thickness dependence,
whereas prefix--suffix-determined typing gives the
main higher-rank transfer. The finite-abelian-group
extension has its own proof in
Appendix~\ref{app:window-abelian}.
The companion CFG result~\cite{kuriyamaCFG}
separately excludes a universal original-thickness
polynomial for arbitrary finite typing at fan-out one;
the semilattice typing of that example is
not determined by any fixed prefix--suffix window.
Polynomial update time in the observed sample is independent of these
data-length refinements
(Theorem~\ref{thm:poly} and Corollary~\ref{cor:capped-gold-poly}).

\subsection{Thickness and reusable canonical contexts}
\label{subsec:thickness-contexts}

\paragraph{Rule thickness and Yoshinaka's efficiency criterion.}
For the quantitative results, presentations are reduced and \emph{good}, as in
Yoshinaka's fixed-dimension, fixed-rank good form
\cite[Section~2.3, after Example~2]{Yoshinaka2011}.  For a rule $R$ of $G$ let
$\theta_G(R)$ be the minimum length of a successful sentence derivation using
$R$, put $\theta_G:=\max_{R\in P}\theta_G(R)$, and define
\[
 t_G:=\max_A\min\{\|\tuple w\|_1:\tuple w\in L_A(G)\}.
\]
Thus $\theta_G$ is rule thickness and $t_G$ nonterminal thickness.
Yoshinaka writes these quantities as $\tau_G$ and $t_G$, respectively, with
$\tau_G\le\|G\|t_G$ in his notation
\cite[Section~4.3]{Yoshinaka2011}.

\begin{lemma}[Polynomial equivalence of the two thickness parameters]
\label{lem:thickness-equivalence}
For every reduced good presentation $G$,
\[
  t_G\le\theta_G
  \qquad\text{and}\qquad
  \theta_G\le 2\|G\|\max\{1,t_G\}.
\]
\end{lemma}

\begin{proof}
For $t_G\le\theta_G$, choose any rule with left-hand side $A$ and place it
in a successful derivation.  Nondeletion preserves the tuple produced at that
occurrence of $A$, so its total length is bounded by the final sentence.
Minimize for $A$ and maximize over $A$.

For the converse, fix a rule $R$ with left-hand side $A$ and follow a shortest
dependency path from the start symbol to $A$.  The path repeats no
nonterminal.  Keep the path child open, complete each sibling by a shortest
terminal tuple, then apply $R$ and complete its children likewise.  Across the
distinct rules on this construction, both explicit terminal material and
child-variable occurrences are bounded by $\|G\|$, giving a successful
sentence of length at most
$\|G\|+\|G\|t_G\le2\|G\|\max\{1,t_G\}$.
\end{proof}

Hence the two thickness parameters are polynomially equivalent on reduced good
presentations.  Yoshinaka's efficiency criterion asks for polynomial
hypothesis construction, characteristic-set cardinality, and total positive
data in target size and rule thickness
\cite[Definition~1]{Yoshinaka2011}; see also
\cite{EyraudHeinzYoshinaka2016}.  Here characteristic-data bounds concern the
capped set-driven reconstruction, while TxtEx stabilization is supplied by the
conservative learner.  Proposition~\ref{prop:structural-good-thickness}
transfers the size--thickness scale for reduced $\lambda$-free inputs; genuine
empty components may instead cause the exponential growth of
Proposition~\ref{prop:empty-component-blowup}.

\begin{lemma}[Reusable canonical-context bound]

\label{lem:reusable-context-bound}
Let $\mathcal G_{\mathrm t}$ be a trimmed good typed MCFG
with $n_{\mathrm{typ}}$ non-start typed nonterminals and maximum non-start rule rank at most
$r_0\ge1$.  Suppose that every non-start typed nonterminal $Y$ has a canonical tuple
$\omega(Y)\in L_Y(\mathcal G_{\mathrm t})$ with
$\|\omega(Y)\|_1\le b_{\mathrm{can}}$, and that every non-start typed rule writes at
most $b_{\mathrm{rule}}$ terminal symbols explicitly.  Then the
canonical contexts may be chosen so that
\[
  |\chi(X)|_{\mathrm{fix}}
  \le n_{\mathrm{typ}}\bigl(b_{\mathrm{rule}}+(r_0-1)b_{\mathrm{can}}\bigr)
\]
for every non-start typed nonterminal $X$.  With these contexts and canonical tuples,
every selected witness word has length at most
\[
  n_{\mathrm{typ}}\bigl(b_{\mathrm{rule}}+(r_0-1)b_{\mathrm{can}}\bigr)+b_{\mathrm{rule}}+r_0b_{\mathrm{can}}.
\]
\end{lemma}

\begin{proof}
If $\widetilde S\to X$ is present take $\chi(X)=\blank_1$; otherwise follow a
shortest typed-nonterminal dependency path to $X$.  It has at most
$n_{\mathrm{typ}}$ non-start steps.  Keeping the path child open and filling
each sibling canonically costs at most
$b_{\mathrm{rule}}+(r_0-1)b_{\mathrm{can}}$ terminals per step.  Goodness
preserves the selected components, their order, and nonempty separators,
giving the context bound.  Adding either one canonical tuple or one rule
template with at most $r_0$ child tuples gives the displayed witness bound.
\end{proof}

\subsection{Refined-Thickness Baseline}
\label{subsec:typed-thickness}

For a reduced good presentation $G$ of an $\eps$-free target, let
$\widetilde G_0$ be its trimmed full yield-type refinement and put
\[
 \bar\theta_{\mathrm{typ}}
 :=\max\{1,\theta_{\widetilde G_0}\},
\]
where $\theta_{\widetilde G_0}$ is rule thickness computed in the refined
grammar.  No comparison with the original untyped thickness is assumed.

\begin{lemma}[Refined-thickness baseline]
\label{lem:typed-thickness-baseline}
Fix $f\ge1$, $r_{\max}\ge1$, a finite alphabet $\Sigma$, and a finite
typing $h:\Sigma^*\to M$.  Let $G$ be a reduced good fan-out-$\le f$,
rank-$\le r_{\max}$ presentation of an $\eps$-free language in
$\Cmcf{f}{h}$, and let $\widetilde G_0$ be its trimmed yield-type refinement.
Then the target witness set has cardinality polynomial in $\|G\|$ and total
positive length polynomial in $\|G\|$ and
$\bar\theta_{\mathrm{typ}}:=\max\{1,\theta_{\widetilde G_0}\}$.
\end{lemma}

\begin{proof}
The trimmed refinement $\widetilde G_0$ remains reduced and good.  A
minimum-yield successful derivation using a typed rule with left-hand side
$X$ contains an $X$-tuple of total length at most
$\bar\theta_{\mathrm{typ}}$, so use such tuples canonically.  If
$\widetilde S\to X$ exists take $\chi(X)=\blank_1$; otherwise remove the
chosen $X$-subtree from the same derivation to obtain a canonical context with
fixed material at most $\bar\theta_{\mathrm{typ}}$.  Each typed rule writes at
most that many terminals and has at most $r_{\max}$ canonical children, so
every selected witness has length at most
$(r_{\max}+2)\bar\theta_{\mathrm{typ}}$.  Therefore
\[
 |\Wit(\widetilde G_0)|
 \le |V||M|^f+|P||M|^{fr_{\max}},
 \qquad
 \|\Wit(\widetilde G_0)\|_+
 \le
 (|V||M|^f+|P||M|^{fr_{\max}})
 (1+(r_{\max}+2)\bar\theta_{\mathrm{typ}}),
\]
which is polynomial for fixed $f,r_{\max},h$.
\end{proof}

For trivial typing, no nonterminal is split by
yield type; consequently the original
thickness scale is retained. This recovers
Yoshinaka's untyped characteristic-data
estimate~\cite[Lemma~7]{Yoshinaka2011}.
Without this uniqueness, the refinement
may have much larger thickness, as the next
proposition demonstrates.

\begin{proposition}[Exponential blow-up of full yield-type refinement]
\label{prop:typed-thickness-exponential}
There are a fixed two-element typing $h_{\mathtt c}$, reduced good fan-out-one rank-two
grammars $G_n$ of size $O(n)$, and a constant bound
$\theta_{G_n}\le2$, such that $L(G_n)=\{\mathtt a,\mathtt c\}^+$ and the
trimmed full yield-type refinement $\widetilde G_{n,0}$ satisfies
\[
  \theta_{\widetilde G_{n,0}}\ge 2^n.
\]
\end{proposition}

\begin{proof}
Let $M_{\mathtt c}=\{1,\zeta\}$ with $\zeta^2=\zeta$,
$h_{\mathtt c}(\mathtt a)=1$, and
$h_{\mathtt c}(\mathtt c)=\zeta$.  For $n\ge1$, view
\[
\begin{aligned}
 S&\to SS\mid\mathtt a\mid\mathtt c\mid A_0\mid\cdots\mid A_n,\\
 A_0&\to\mathtt a,\qquad
 A_i\to A_{i-1}A_{i-1}\mid\mathtt c\quad(1\le i\le n)
\end{aligned}
\]
as a fan-out-one MCFG.  It is reduced and good, has size $O(n)$ and rule rank
two, generates $\{\mathtt a,\mathtt c\}^+$, and every rule has a successful
use of length at most two, so $\theta_{G_n}\le2$.

In the full refinement, the retained copy $(A_n)_1$ can derive only words
without $\mathtt c$; recursively the unique such yield is
$\mathtt a^{2^n}$.  Hence the typed rule
$S_{(1)}\to(A_n)_1$ has successful thickness at least $2^n$.
\end{proof}

Thus full yield-type refinement admits no bound polynomial in
$\|G\|\max\{1,\theta_G\}$ for arbitrary fixed typing, already at fan-out one
and rank two.  This is a limitation of the full-refinement route, not of the
language $\{\mathtt a,\mathtt c\}^+$ itself.

\paragraph{Prior lower bound at fan-out one.}
The revised companion CFG paper proves an
operator-independent exponential characteristic-data
lower bound for fixed two-element semilattice typing,
already at fan-out one and rule rank two
\cite{kuriyamaCFG}.
We do not reproduce this prior CFG result.
It does not automatically extend to higher fan-out,
whose tuple substitutability tests more arities.

\subsection{Rule Rank One: Grammar-Size Polynomial Data}
\label{subsec:rankone-polydata}

Unary derivation chains admit a grammar-size bound
independent of thickness, even under arbitrary fixed
componentwise typing.

\begin{theorem}[Rank-one polynomial construction and characteristic data]
\label{thm:rankone-polydata}
Fix $f\ge1$, finite $\Sigma$, and a typing
$h:\Sigma^*\to M$ into a finite monoid.
For any finite standard fan-out-$\le f$,
rank-$\le1$ grammar $G_{\mathrm{in}}$
of $L\in\Cmcf{f}{h}$,
$\mathcal B_{f,h}^{\le1}$ is polynomial-time in
$\|K\|_+$ and has a characteristic sample of
total length $\operatorname{poly}_{f,h}(\|G_{\mathrm{in}}\|)$.
The direct-start datum $\eps$ is added when
$\eps\in L$ without affecting this bound.
The conservative learner
$\mathcal A_{f,h}^{\le1}$ consequently
TxtEx-identifies
$\Cmcf{f}{h}\cap\mathsf{MCFL}(f,1)$
with polynomial time per datum.
\end{theorem}

\begin{proof}
Normalize without increasing rank
(Proposition~\ref{prop:good-normalization}),
and trim the typed refinement.
Write $n_{\mathrm{typ}}\le|V||M|^f$ for its
number of non-start typed nonterminals and
$b_G\le\max\{1,\|G\|\}$ for the largest explicit
terminal count of a rule, taking $b_G\ge1$.
Choose a typed terminal unary derivation that
minimizes first total yield length and then the
number of rules. No typed nonterminal can repeat:
splicing in the later suffix would not increase
yield length and would strictly shorten the
chain. Therefore
$\|\omega(X)\|_1\le n_{\mathrm{typ}}b_G$.
Lemma~\ref{lem:reusable-context-bound} bounds
each witness word by $(2n_{\mathrm{typ}}+1)b_G$,
while Proposition~\ref{prop:witness-cardinality}
gives at most $n_{\mathrm{typ}}+|P||M|^f$
witnesses. Hence
\[
 \|\Wit(\widetilde G_0)\|_+
 \le(n_{\mathrm{typ}}+|P||M|^f)
        \bigl(1+(2n_{\mathrm{typ}}+1)b_G\bigr).
\]
The polynomial overhead of normalization and the
time and convergence bounds follow from
Theorem~\ref{thm:poly} and
Corollary~\ref{cor:capped-gold-poly}.
\end{proof}

This extends the rank-one untyped scale
\cite[Lemma~7]{Yoshinaka2011} and the fan-out-one
argument of~\cite{kuriyamaCFG} to arbitrary
fixed finite typing at every fixed fan-out.
Rank two cannot generally yield a
grammar-size-only bound: the size-$O(n)$
grammar $A_0\to\mathtt c$,
$A_i\to A_{i-1}A_{i-1}$, $S\to A_n$
generates the singleton $\{\mathtt c^{2^n}\}$,
which is $(f,h)$-tuple-substitutable for
every fixed $f,h$: over a one-letter alphabet,
sharing a sentence context forces equal total
component lengths, so every context accepts
both tuples or neither. Every characteristic
sample for this singleton target contains a word
of length at least $2^n$.

\subsection{接頭辞--接尾辞窓で決まる型付け：厚みに依存する特徴データ}
\label{subsec:boundary-polydata}

補題~\ref{lem:typed-thickness-baseline} は，任意の $h$ に対する特徴データ上界を，
型付き精密化文法の厚みによって与える．
命題~\ref{prop:typed-thickness-exponential} によれば，
この精密化後の厚みが元の良形表示のパラメータで多項式に抑えられるとは限らない．
それでも，有限型付けが固定接頭辞--接尾辞窓によって決まる場合には，
より強い転移結果が得られる．
この境界型付けの場合を，Yoshinaka~\cite[Section~4.3]{Yoshinaka2011} と同様の，
厚みを考慮した文法表示の尺度で定式化する．

$G=(V,\Sigma,P,S,\fan)$ を縮約済み良形 MCFG とし，
\[
  \bar\theta_G:=\max\{1,\theta_G\}
\]
とおく．$\theta_G$ は既に定義した規則厚みである．
以下の議論は $\lambda$-自由な良形の核心部分について述べる．
標的が $\eps$ も含む場合，定理~\ref{thm:exact-reconstruction} と同様に，
$L\setminus\{\eps\}$ に議論を適用してから，
開始記号に証拠 $\eps$ を一つ加えればよい．
これは長さ0の例を一つ増やすだけで，どちらの多項式上界にも影響しない．

\begin{lemma}[規則厚みによる局所的な終端記号数の上界]
\label{lem:thickness-local-bounds}
$G$ を縮約済み良形文法とする．このとき：
\begin{enumerate}[label=(\roman*),leftmargin=*]
\item 任意の規則 $R$ のテンプレートに明示的に書かれた
終端記号の総数は $\theta_G$ 以下である．
\item 任意の非終端記号 $A$ に対し，
$\|\tuple u\|_1\le\theta_G$ を満たす
終端タプル $\tuple u\in L_A(G)$ が存在する．
\end{enumerate}
\end{lemma}

\begin{proof}
縮約性により，すべての規則は何らかの成功導出に出現する．
良形文法は非消去なので，その規則が明示的に書く終端記号と，
左辺で生成されるタプルの全成分は，最終的な文の中にそれぞれちょうど一度残る．
したがって，その規則を用いる最短の成功導出の生成語長は，
規則の明示的終端記号数と，その出現箇所の左辺タプルの総長の両方を上から抑える．
規則全体について最大を取ると (i), (ii) が従う．
\end{proof}

\begin{lemma}[型なし規則厚みによる正準文脈の制御]
\label{lem:untyped-contexts}
$G=(V,\Sigma,P,S,\fan)$ を縮約済み良形 MCFG とし，
$\widetilde G_0$ を任意の固定有限型付け $h$ による生成語型付き精密化の
刈り込み結果とする．
ファンアウト $d$ の各 $A\in V$ に対し，
次を満たす許容文脈 $E_A\in\mathcal E_d$ が存在する：
\[
 |E_A|_{\mathrm{fix}}\le\bar\theta_G,
 \qquad E_A[\tuple u]\in L(G)
 \quad\text{for all }\tuple u\in L_A(G).
\]
さらに $E_S=\blank_1$ と選んでよい．
したがって，開始記号以外の型付き非終端記号 $X=A_{\mathbf t}$ に対し，
$\chi(X):=E_A$ は定義~\ref{def:target-canonical} を満たす．
各正準タプルについて
$\|\omega(X)\|_1\le b_{\mathrm{can}}$ が成立し，
最大規則ランクを $r_{\max}$ とすると，
$\Wit(\widetilde G_0)$ に属するすべての語の長さは
\[
  2\bar\theta_G+r_{\max}b_{\mathrm{can}}
\]
以下である．
\end{lemma}

\begin{proof}
$E_S=\blank_1$ とする．
$A\ne S$ に対しては，左辺が $A$ である元文法の規則 $R$ を一つ選ぶ．
$G$ は縮約済みなので，$R$ を使用し，生成語長が
$\theta_G(R)\le\theta_G$ である成功導出が存在する．
その導出中の $A$ の出現を根とする部分木を取り除き，
その $d$ 成分を順序付きの穴で置き換える．
補題~\ref{lem:filling-identity} を根への経路に沿って繰り返し適用すると，
良形性によって得られる文脈 $E_A$ は $\mathcal E_d$ に属する．
その固定終端部分は選んだ文の一部なので，長さは
$\theta_G(R)$ 以下である．
$L_A(G)$ の任意のタプルをこの位置へ代入すると，
元文法における別の成功導出が得られる．
よって $E_A$ は $L_A(G)$ のすべてのタプルを受理する．

命題~\ref{prop:output-invariants}(i) より
$L_X\subseteq L_A(G)$ だから，
$E_A$ は型付き記号 $X$ のすべてのタプルも受理する．
この周囲の型なし導出は，
命題~\ref{prop:output-invariants}(ii) によって
完全な型付き精密化へ持ち上げることができる．
固定された兄弟部分木の導出の下では $X$ の型から祖先の型が定まり，
新しい開始記号はどの開始型でも許すからである．
したがって $\chi(X)$ は必要な共通文脈の性質を満たす．
開始記号からの直接リンクがある場合には，
$X=S_{\mathbf t}$ であり，$\chi(X)=\blank_1$ とすればよい．

証拠語は
$\chi(X)[\omega(X)]$，
$\chi(X)[\tuple c]$，または
$\chi(X)[\rho(\omega(Y_1),\ldots,\omega(Y_r))]$
のいずれかである．
文脈の終端部分の寄与は高々 $\bar\theta_G$ であり，
元の規則が明示的に書く終端記号数も
補題~\ref{lem:thickness-local-bounds}(i) により
高々 $\bar\theta_G$ である．
さらに子のタプル数は高々 $r_{\max}$ なので，
ランク0の場合も含めて表示した上界を得る．
\end{proof}

$k,\ell\ge0$ を固定し，
$h\preceq h_\Sigma^{k,\ell}$ を満たす
任意の有限型付け $h:\Sigma^*\to M$ を考える．
すなわち $h$ は，全射な接頭辞--接尾辞型付け
$h_\Sigma^{k,\ell}$ を通じて因子化する．
$G$ のファンアウトは高々 $f$，
規則ランクは固定された $r_{\max}\ge1$ 以下とする．
$\widetilde G_0$ は $G$ の完全 $h$-型付き精密化の刈り込み結果とし，
残存する開始記号以外の $h$-型付き非終端記号の数を
$n_{\mathrm{typ}}$ とする．

\begin{lemma}[短い接頭辞--接尾辞型付き正準タプル]
\label{lem:boundary-short-canonical}
$\widetilde G_0$ の開始記号以外の任意の型付き非終端記号 $X$ に対し，
正準タプルを
\[
  \|\omega(X)\|_1
  \le f(k+\ell+1)n_{\mathrm{typ}}(r_{\max}+1)\bar\theta_G
\]
となるように選べる．
\end{lemma}

\begin{proof}
刈り込んだ型付き文法 $\widetilde G_0$ で，
$X$ からの終端導出のうち，
文法規則の節点数が最小のものを選ぶ．
次に，根の各成分の $h_\Sigma^{k,\ell}$-型を決定する
終端位置を標識する．
定義~\ref{def:prefix-suffix-typing} の
$m_0=\max\{1,k+\ell\}$ より短い成分ではすべての終端位置を標識し，
それ以外の成分では先頭 $k$ 個と末尾 $\ell$ 個を標識する．
$k=\ell=0$ のときは任意の終端位置を一つ標識する．
標識された終端位置は合計高々 $f(k+\ell+1)$ 個である．

最小節点数のこの導出では，
根から葉への経路に同じ型付き非終端記号が二度出現することはない．
もし上側とその真の子孫に同じ型付き標識があれば，
下側の部分木は上側と同じ型付き非終端記号からの
正しい終端導出である．
上側を下側の部分木で置き換えれば，
規則節点数は真に減るので矛盾する．
この導出に出現する型付き非終端記号は
すべて $\widetilde G_0$ に残存している．
したがって，根から標識終端位置への各経路の
文法節点数は高々 $n_{\mathrm{typ}}$ である．
それらの経路の合併に含まれる節点数は高々
$f(k+\ell+1)n_{\mathrm{typ}}$ となる．

標識された根から葉への経路の合併を，保持する部分導出の骨格とする．
線形性と非消去性により，
取り除いた各子のそれぞれの成分は，
根の一つの成分の中の連続区間に寄与する．
一つの子成分が根の複数成分へ分割されることはない．
骨格から除外した各子部分木を，
対応する元の型なし非終端記号が生成する最短の終端タプルで置換する．
補題~\ref{lem:thickness-local-bounds}(ii) により
このタプルの総文字数は $\bar\theta_G$ 以下であり，
同 (i) により骨格に残した各規則の
明示的終端記号数も $\bar\theta_G$ 以下である．
良形性によって置換後の各成分は非空である．

骨格から除外した子は，
短い根成分には寄与し得ない．
短い成分のすべての終端位置は標識されているため，
子が非空の区間に寄与するならば，
そこに標識された終端位置が含まれ，
その位置への経路が子を骨格に残すからである．
長い根成分においては，
置換されるブロックは保護された接頭辞位置や接尾辞位置を横切らない．
したがって，置換は各成分の標識されていない内部だけで起こり，
標識した先頭 $k$ 個と末尾 $\ell$ 個の終端記号がそのまま保存される．
また，その成分は置換後も長い語の側に留まる．
よって根の接頭辞--接尾辞型は不変である．
$k=\ell=0$ では唯一の標識位置が非空性を保存する．
$h\preceq h_\Sigma^{k,\ell}$ の因子写像を適用すれば，
根の成分別 $h$-型も変化しない．

骨格の各節点が明示的に書く終端記号数は高々
$\bar\theta_G$，取り除いた子部分木の数は高々
$r_{\max}$ である．
したがって，新しい根タプルの長さは高々
\[
  f(k+\ell+1)n_{\mathrm{typ}}(r_{\max}+1)\bar\theta_G
\]
となる．
このタプルは $X$ を添字づける $h$-型をもつ．
置換によって元の型付き導出の内部標識が変わることはあるが，
これは問題にならない．
元文法 $G$ 上の導出と保存された根の型だけを残し，
命題~\ref{prop:output-invariants}(ii) により
その導出を根 $X$ をもつ $G^h$ の導出へ一意に持ち上げる．
$X$ は $\widetilde G_0$ に残る記号なので，
持ち上げた導出を $X$ を通る成功導出へ挿入でき，
刈り込み後にもこの導出が残る．
その根のタプルを $\omega(X)$ として選べる．
\end{proof}

補題~\ref{lem:untyped-contexts} を
\[
 b_{\mathrm{can}}
 =f(k+\ell+1)n_{\mathrm{typ}}
       (r_{\max}+1)\bar\theta_G
\]
として適用する．
すると各証拠語の長さは高々
$2\bar\theta_G+r_{\max}b_{\mathrm{can}}
 =O_{f,r_{\max},k,\ell,h,\Sigma}
   (\|G\|\bar\theta_G)$
となる．ここで
$n_{\mathrm{typ}}\le|V||M|^f$ を用いた．

\begin{theorem}[接頭辞--接尾辞窓で決まる型付けの多項式特徴データ]
\label{thm:boundary-polydata}
$f\ge1$，$r_{\max}\ge1$，$k$，$\ell$，
有限アルファベット $\Sigma$，
および $h\preceq h_\Sigma^{k,\ell}$ を満たす
有限型付け $h:\Sigma^*\to M$ を固定する．
$G$ を $\Cmcf{f}{h}$ に属する言語の
縮約済み良形表示とし，そのファンアウトは高々 $f$，
規則ランクは高々 $r_{\max}$ とする．
このとき
$\mathcal B_{f,h}^{\le r_{\max}}$ は正例標本のサイズについて
多項式時間で計算できる．
さらに，特徴標本を
\[
 |\Wit(\widetilde G_0)|=O_{f,r_{\max},h}(\|G\|),
 \qquad
 \|\Wit(\widetilde G_0)\|_+
 =O_{f,r_{\max},k,\ell,h,\Sigma}
   (\|G\|^2\bar\theta_G)
\]
を満たすように選べる．
ここで $\bar\theta_G=\max\{1,\theta_G\}$ である．
消去規則を許した縮約済み $\lambda$-自由入力表示
$G_{\mathrm{in}}$ に対しても，
命題~\ref{prop:structural-good-thickness} の
構造的良形変換を通せば，
特徴データの総長は
$\|G_{\mathrm{in}}\|\max\{1,t_{G_{\mathrm{in}}}\}$
の多項式で抑えられる．
\end{theorem}

\begin{proof}
$h$-精密化文法の型付き非終端記号数は高々
$|V||M|^f$，
型付き規則の個数は高々
$|P||M|^{fr_{\max}}$ である．
より細かい $h_\Sigma^{k,\ell}$ の境界型を保存すれば
それを通じて決まるより粗い $h$-型も保存されるため，
補題~\ref{lem:boundary-short-canonical} が適用できる．
さらに補題~\ref{lem:untyped-contexts} から，
各証拠語の長さは
\[
 O_{f,r_{\max},k,\ell,h,\Sigma}
 (\|G\|\bar\theta_G)
\]
以下となる．
命題~\ref{prop:witness-cardinality} により
\[
 |\Wit(\widetilde G_0)|
 \le n_{\mathrm{typ}}+|P||M|^{f r_{\max}},
\]
ただし $n_{\mathrm{typ}}\le|V||M|^f$ である．
証拠語の多項式個数に上の各語の長さ上界を掛け，
必要なら $\eps$ の証拠を高々一つ加えると，
主張した正例総文字数の上界を得る．
同じ成功導出に基づく証拠は
型付き精密化文法の規則厚みにも多項式上界を与えるが，
この中間パラメータは上界の証明には不要である．
仮説文法の構成時間に関する主張は
定理~\ref{thm:poly} に従う．
\end{proof}

入力表示に関する主張は
命題~\ref{prop:structural-good-thickness} による．
消去規則を含む入力文法では，
良形変換に先立って終端記号が消えることがあり得るため，
非終端記号の厚みを用いる．
系~\ref{cor:capped-gold-poly} は，
対応するランク制限付き保守的学習器の
各データ処理における多項式時間上界を与える．

ファンアウト1では，この窓型付け上界自体は新しくない．
Yoshinaka の $(k,\ell)$-代入可能な CFG の学習器は，
すでに文法サイズと厚みに関する多項式特徴データを与える
~\cite[Theorem~1]{Yoshinaka2008}．
また，関連する改訂 CFG 論文は，窓型付けを有限群値で精密化した場合へ
結果を拡張している~\cite{kuriyamaCFG}．
補題~\ref{lem:boundary-short-canonical} の寄与は，
\emph{複数の}タプル成分の境界情報を同時に保存し，
任意の固定ファンアウト $f\ge2$ でも同種の上界を実現する点にある．

上界は縮約済み良形の規則ランク有界な文法表示に相対的なものである．
後述の定理~\ref{thm:intrinsic-ranktwo-witness} は，
この高ランクの場合が言語レベルで必要であることを示す．
すなわち，競合する文法が任意の有限ファンアウトを使ってよい場合でも，
最小ファンアウトと最小規則ランクがともに2である，
$(1,1)$-境界型付けされた非文脈自由標的を与える．
後の例では，$L_{\mathrm{intr}}$ が窓型付けの定理の対象となり，
ランク1の言語 $L_{\mathrm{cd}}$ と $L_\star$ は
定理~\ref{thm:rankone-polydata} の対象となる．
したがって，いずれの分離結果も
対応する定量的学習上界と両立している．

\section{多次元代入可能性との比較}
\label{sec:comparison}

Yoshinaka's multidimensional substitutability is the direct positive-data
baseline~\cite{Yoshinaka2009,Yoshinaka2011}.  Under trivial typing the present
framework reproduces that semantic class; finite typing refines which observed
tuples may be identified while preserving compositional rule witnesses.
The MCFG proofs additionally require componentwise type propagation and, for
the quantitative boundary theorem, simultaneous control of several tuple
components.  Yoshinaka's $(k,\ell)$ hierarchy is a separate CFG result
\cite{Yoshinaka2008}; $h_\Sigma^{k,\ell}$ lifts the same boundary
information componentwise.  Since every $(f,h)$ condition includes arity one,
\[
  \Cmcf{f}{h}\cap\mathrm{CFL}\subseteq\Cmcf{1}{h}.
\]
The examples below separate finite typing from the trivial baseline,
prefix--suffix typings, and regular-filter explanations.

\subsection{Yoshinaka's Baseline and Prefix--Suffix Specialization}

\begin{definition}[Yoshinaka multidimensional substitutability]
\label{def:yoshinaka-pd-substitutability}
Using the sentence-context sets $\mathcal E_d$ defined in
Section~\ref{sec:prelim}, a language $L$ is \emph{$f$D-substitutable} if,
for every $1\le d\le f$, all $E,F\in\mathcal E_d$, and all
$\tuple x,\tuple y\in(\Sigma^+)^d$,
\[
 E[\tuple x],\quad E[\tuple y],\quad F[\tuple x]\in L
 \Longrightarrow
 F[\tuple y]\in L.
\]
Write $\YSub{f}$ for this semantic class and
$\YSL{f}{r}:=\YSub{f}\cap\mathsf{MCFL}(f,r)$, where $r$ bounds rule rank.
We refer to the displayed four-instance implication as the
\emph{substitution implication}.
Yoshinaka's original definition quantifies over every arity $d\le f$, uses
nonempty tuple components, and requires nonempty internal separators in
$d$-contexts; footnote~3 separately discusses the stronger alternative that
permits empty components and separators
\cite[Section~3.1 and footnote~3]{Yoshinaka2011}.
\end{definition}

Recall the finite monoid morphism $h_\Sigma^{k,\ell}$ from
Definition~\ref{def:prefix-suffix-typing}.  In Yoshinaka's word-level
definition, $w_1,y_1,w_1',w_2,y_2,w_2'\in\Sigma^*$,
$u\in\Sigma^k$, and $v\in\Sigma^\ell$, with
$uy_1v,uy_2v\ne\eps$; the condition is
\[
 w_1uy_1vw_1',\ w_1uy_2vw_1'\in L
 \quad\Longrightarrow\quad
 \bigl(w_2uy_1vw_2'\in L\iff w_2uy_2vw_2'\in L\bigr)
\]
\cite{Yoshinaka2008}.  Thus the compared fragments $uy_iv$ are nonempty even
though the middle words $y_i$ themselves may be empty.  Equality of their
$h_\Sigma^{k,\ell}$-values is exactly equality of the boundary information
used in this condition, including the exact identity of words shorter than
$m_0=\max\{1,k+\ell\}$.

\paragraph{Specializations of the semantic condition.}
The trivial typing removes the componentwise
type premise, so
\[
 \Cmcf{f}{h_{\mathrm{triv}}}
 =\YSub{f}\cap f\text{-}\mathrm{MCFL}
 \subseteq\Cmcf{f}{h}.
\]
At fan-out one,
$(1,h_\Sigma^{k,\ell})$-substitutability
is precisely Yoshinaka's word-level
$(k,\ell)$ condition~\cite{Yoshinaka2008}:
equal window types mean identical short words
or common prefixes and suffixes of lengths
$k$ and $\ell$, respectively.
The same correspondence applies
componentwise at larger fan-out.

\paragraph{Compatibility with Yoshinaka's efficiency scale.}
For trivial typing no productive nonterminal
has more than one yield type. Yoshinaka's
untyped grammar witnesses already satisfy
$|K_G|\le|P|$ and
$\|K_G\|\le|P|\tau_G$
\cite[Lemma~7]{Yoshinaka2011}, with polynomial
time and data for fixed dimension and rank
\cite[Theorem~1]{Yoshinaka2011}.
Thus the present quantitative claims concern
new typed, multi-component settings rather
than re-proving the trivial-typing result.

\subsection{An Intrinsically Rank-Two Boundary-Typed Target}
\label{subsec:intrinsic-ranktwo}

The next example makes the distinction between the rank-one and higher-rank
quantitative regimes intrinsic to the language.  Its rule-rank obstruction is
not created by finite typing: the balanced-count factor used below already
belongs to Yoshinaka's trivial-typing class.  The construction instead combines
that obstruction with minimum fan-out two and a separation from the
trivial-typing class inside one boundary-typed target.

Let
\[
 \Sigma_{\mathrm{intr}}
 :=\{\mathtt a,\mathtt b,\mathtt l,\mathtt r,\mathtt p,\mathtt q,
       \mathtt c,\mathtt d,\mathtt{\#},\mathtt{\$},
       \mathtt u,\mathtt v,\mathtt w\}
\]
and put
\[
 P_{\mathrm{intr}}
 :=\{\mathtt{lpc},\mathtt{lpd},\mathtt{lqc},
       \mathtt{rpc},\mathtt{rpd}\}.
\]
Define
\[
 \mathrm{MIX}_2^+
 :=\{x\in\{\mathtt a,\mathtt b\}^+:
       |x|_{\mathtt a}=|x|_{\mathtt b}\},
 \qquad
 T_3^+
 :=\{\mathtt u^n\mathtt v^n\mathtt w^n:n\ge1\},
\]
and
\[
 L_{\mathrm{intr}}
 :=P_{\mathrm{intr}}\,\mathtt{\#}\,\mathrm{MIX}_2^+\,
   \mathtt{\$}\,T_3^+.
\]

To make the role of the witness explicit, let
\[
 \varphi_{\mathrm{mix}}:\{\mathtt a,\mathtt b\}^*\to\mathbb Z,
 \qquad
 \varphi_{\mathrm{mix}}(\mathtt a)=1,\quad
 \varphi_{\mathrm{mix}}(\mathtt b)=-1.
\]
Then
\[
 \mathrm{MIX}_2^+
 =\{\mathtt a,\mathtt b\}^+\cap\varphi_{\mathrm{mix}}^{-1}(0).
\]
The language $\{\mathtt a,\mathtt b\}^+$ is trivially
$f$-tuple-substitutable, so
Proposition~\ref{prop:abelian-coset-filter}(ii) gives
$\mathrm{MIX}_2^+\in\YSub{f}$ for every $f\ge1$; the rank-two CFG for $B$
displayed below also shows that this ingredient is context-free.  Thus the
intrinsic rank-two phenomenon used here already occurs inside the
trivial-typing baseline.  The role of $L_{\mathrm{intr}}$ is to combine that
known rule-rank obstruction with non-context-freeness, minimum fan-out two,
and a separation from $\YSub{f}$ in a single finitely typed target.

\begin{theorem}[A non-context-free intrinsic-rank-two typed witness]
\label{thm:intrinsic-ranktwo-witness}
The language $L_{\mathrm{intr}}$ satisfies
\[
 L_{\mathrm{intr}}\in\mathsf{MCFL}(2,2)\setminus\mathrm{CFL},
\]
and, for every $f\ge2$,
\[
 L_{\mathrm{intr}}
 \in\Cmcf{f}{h_{\Sigma_{\mathrm{intr}}}^{1,1}}
 \setminus\YSub{f}.
\]
Moreover, for every finite $f'\ge1$,
\[
 L_{\mathrm{intr}}\notin\mathsf{MCFL}(f',1).
\]
Consequently the minimum fan-out of $L_{\mathrm{intr}}$ is exactly two, and
its minimum rule rank among all finite-fan-out MCFG presentations is exactly
two.  In particular, Theorem~\ref{thm:boundary-polydata} applies to a
non-context-free language that cannot be covered by
Theorem~\ref{thm:rankone-polydata} at any finite fan-out.
\end{theorem}

\begin{proof}
First we give a reduced good fan-out-two grammar of maximum rule rank two.
Let $B$ generate the nonempty balanced-count language by
\[
\begin{aligned}
 B\to{}&
 \mathtt a B\mathtt b B
 \mid \mathtt b B\mathtt a B
 \mid \mathtt a B\mathtt b
 \mid \mathtt a\mathtt b B
 \mid \mathtt a\mathtt b\\
 &{}\mid
 \mathtt b B\mathtt a
 \mid \mathtt b\mathtt a B
 \mid \mathtt b\mathtt a .
\end{aligned}
\]
The usual first-return induction on
$|x|_{\mathtt a}-|x|_{\mathtt b}$ gives
$L_B=\mathrm{MIX}_2^+$.

Let the fan-out-two nonterminal $A$ have rules
\[
 A\to(\mathtt u,\mathtt w)
 \mid
 (\mathtt u\xi_1^1\mathtt v,\mathtt w\xi_1^2)(A).
\]
It derives exactly the tuples
\[
 \{(\mathtt u^n\mathtt v^{n-1},\mathtt w^n):n\ge1\}.
\]
For every $\alpha\in P_{\mathrm{intr}}$, add the start rule
\[
 S\to
 \alpha\mathtt{\#}\xi_1^1\mathtt{\$}
 \xi_2^1\mathtt v\xi_2^2(B,A).
\]
The resulting grammar generates exactly $L_{\mathrm{intr}}$.  It is reduced,
$\lambda$-free, linear, nondeleting, and nonpermuting.  It is also nonmerging:
the two components of $A$ are separated by the literal terminal $\mathtt v$
in every start rule.  Hence it is a good fan-out-two grammar of maximum rule
rank two.

The language is not context-free.  Intersect it with the regular language
\[
 R:=
 \mathtt{lpc\#ab\$}\,
 \mathtt u^+\mathtt v^+\mathtt w^+ .
\]
Then
\[
 L_{\mathrm{intr}}\cap R
 =
 \{\mathtt{lpc\#ab\$}\mathtt u^n\mathtt v^n\mathtt w^n:n\ge1\}.
\]
A homomorphism erasing the fixed prefix and fixing
$\mathtt u,\mathtt v,\mathtt w$ maps this language onto
$\{\mathtt u^n\mathtt v^n\mathtt w^n:n\ge1\}$, which is not context-free.
Closure of context-free languages under intersection with regular languages
and homomorphism therefore shows that $L_{\mathrm{intr}}$ is not
context-free.  Since fan-out-one MCFGs are exactly context-free grammars in
the present syntax, while the displayed grammar has fan-out two, the minimum
fan-out is exactly two.

For typed membership, let
\[
 Q_{\mathrm{intr}}
 :=P_{\mathrm{intr}}\,\mathtt{\#}\,\{\mathtt a,\mathtt b\}^+\,
   \mathtt{\$}\,\mathtt u^+\mathtt v^+\mathtt w^+ .
\]
This language is strictly $2$-local.  One complete set of allowed adjacent
pairs, including end markers, is
\[
\begin{gathered}
 \vdash\mathtt l,\ \vdash\mathtt r,\quad
 \mathtt{lp},\ \mathtt{lq},\ \mathtt{rp},\quad
 \mathtt{pc},\ \mathtt{pd},\ \mathtt{qc},\\
 \mathtt{c\#},\ \mathtt{d\#},\quad
 \mathtt{\#a},\ \mathtt{\#b},\quad
 \mathtt{aa},\ \mathtt{ab},\ \mathtt{ba},\ \mathtt{bb},\\
 \mathtt{a\$},\ \mathtt{b\$},\quad
 \mathtt{\$u},\quad
 \mathtt{uu},\ \mathtt{uv},\ \mathtt{vv},\ \mathtt{vw},\
 \mathtt{ww},\quad
 \mathtt w\dashv .
\end{gathered}
\]
Define the homomorphism
\[
 \Phi_{\mathrm{intr}}:\Sigma_{\mathrm{intr}}^*\to\mathbb Z^3
\]
by
\[
\begin{aligned}
 \Phi_{\mathrm{intr}}(\mathtt a)&=(1,0,0),&
 \Phi_{\mathrm{intr}}(\mathtt b)&=(-1,0,0),\\
 \Phi_{\mathrm{intr}}(\mathtt u)&=(0,1,0),&
 \Phi_{\mathrm{intr}}(\mathtt v)&=(0,-1,1),&
 \Phi_{\mathrm{intr}}(\mathtt w)&=(0,0,-1),
\end{aligned}
\]
and let every other terminal have image $(0,0,0)$.  Then
\[
 L_{\mathrm{intr}}
 =Q_{\mathrm{intr}}\cap\Phi_{\mathrm{intr}}^{-1}(0).
\]
Indeed, the three coordinates impose respectively
\[
 |x|_{\mathtt a}=|x|_{\mathtt b},\qquad
 |x|_{\mathtt u}=|x|_{\mathtt v},\qquad
 |x|_{\mathtt v}=|x|_{\mathtt w}.
\]
Lemma~\ref{lem:sl2-boundary-control} and
Proposition~\ref{prop:abelian-coset-filter}(ii) show that the displayed
intersection is $(f,h_{\Sigma_{\mathrm{intr}}}^{1,1})$-tuple-substitutable
for every $f\ge1$.  Together with the fan-out-two grammar above, this gives
\[
 L_{\mathrm{intr}}
 \in\Cmcf{f}{h_{\Sigma_{\mathrm{intr}}}^{1,1}}
 \qquad\text{for every }f\ge2.
\]

The language is outside Yoshinaka's trivial-typing class already at arity one.
Take
\[
 x:=\mathtt p,\qquad y:=\mathtt q,
\]
and one-hole contexts
\[
 E:=\mathtt l\blank_1\mathtt{c\#ab\$uvw},
 \qquad
 F:=\mathtt r\blank_1\mathtt{d\#ab\$uvw}.
\]
Then
\[
 E[x],\ E[y],\ F[x]\in L_{\mathrm{intr}},
 \qquad
 F[y]\notin L_{\mathrm{intr}}.
\]
Thus the arity-one substitution implication fails, so
$L_{\mathrm{intr}}\notin\YSub{f}$ for every $f\ge1$.  Notice again that the
boundary typing distinguishes the offending one-letter factors:
\[
 h_{\Sigma_{\mathrm{intr}}}^{1,1}(\mathtt p)
 \ne h_{\Sigma_{\mathrm{intr}}}^{1,1}(\mathtt q).
\]

It remains to prove intrinsic rule rank.  Put
\[
 \mathrm{MIX}_2
 :=\{x\in\{\mathtt a,\mathtt b\}^*:
       |x|_{\mathtt a}=|x|_{\mathtt b}\}.
\]
Gallot calls an MCFG \emph{non-branching} when every positive-rank rule has a
single child nonterminal; in the present terminology this is exactly rule
rank at most one, with no globally fixed fan-out bound
\cite[Definition~38]{Gallot2021}.  Since every finite grammar has a finite
maximum nonterminal arity, Gallot's non-branching MCFL class is
\[
 \bigcup_{f'<\infty}\mathsf{MCFL}(f',1).
\]
Gallot's Theorem~29 identifies this class with the finite-index EDT0L
languages, and Theorem~31 proves that
$\mathrm{MIX}_2$ is not finite-index EDT0L.
Definition~38 explicitly permits a child variable to
occur at most once, and Theorem~28 confirms that
non-permuting non-erasing presentations generate
the same non-branching class; thus the use of
erasing terminal homomorphisms below stays within
the formalism of the cited lower bound.  Equivalently,
$\mathrm{MIX}_2$ is not generated by any non-branching MCFG; Gallot's
Propositions~5--6 give the underlying bounded-$k$-derivation contradiction
\cite[Theorem~29, Theorem~31, Propositions~5--6]{Gallot2021}.

Bishop et al. later give a peer-reviewed expansion of the same obstruction in
the closely corresponding R-MCFG formalism.  Their Definition~6.1 uses
\emph{rank} for nonterminal arity (the present fan-out), while
\emph{non-branching} still means that a propagating rule has a single child
nonterminal.  Lemma~6.8 supplies the normal form; Definition~7.1 identifies
the same balanced language, and Section~7 proves the obstruction, culminating
in the stronger non-EDT0L statement of Theorem~C
\cite[Definition~6.1, Lemma~6.8, Definition~7.1, Section~7,
Theorem~C]{BishopEtAl2026EDT0L}.  This is corroborating evidence; Gallot's
Theorems~29 and~31 remain the direct implication used here.

Now let $\pi$ fix $\mathtt a,\mathtt b$ and erase every other terminal in
$\Sigma_{\mathrm{intr}}$.  Then
\[
 \pi(L_{\mathrm{intr}})=\mathrm{MIX}_2^+.
\]
If $L_{\mathrm{intr}}$ had a fan-out-$f'$ rule-rank-one presentation,
applying $\pi$ to the terminal material in its rule templates would preserve
fan-out and rule rank and would give such a presentation of
$\mathrm{MIX}_2^+$.  Empty terminal pieces created by erasure are permitted
in the standard MCFG formalism.  Adding a fresh start symbol with one unary
identity link to the old start and one rank-zero rule generating $\eps$ would
then give a fan-out-$f'$ rule-rank-one presentation of $\mathrm{MIX}_2$,
contradicting Gallot's non-branching lower bound.  Hence
\[
 L_{\mathrm{intr}}\notin\mathsf{MCFL}(f',1)
 \qquad\text{for every finite }f'\ge1.
\]
Together with the displayed rank-two grammar, this proves the claimed
intrinsic rule rank.
\end{proof}

The construction is deliberately modular: the $\mathrm{MIX}_2^+$ factor
carries the obstruction to rule rank one, while the $T_3^+$ factor forces the
target outside CFL and hence to minimum fan-out two.  Thus the theorem
establishes the coexistence of the two lower bounds in a single finitely typed
target; the two bounds arise from separate factors rather than from an
irreducible interaction between the resources.

The non-branching lower bound itself is prior work.  What is specific to the
present framework is that a strictly local wrapper and the $(1,1)$ boundary
typing place this target inside the typed class, while the displayed four-instance substitution implication keeps it outside Yoshinaka's trivial-typing baseline.
The wrapper retains the prior non-branching obstruction, so the same target
also has intrinsic rule rank two across all finite fan-outs.

\subsection{Finite Typings beyond the Trivial Baseline}
\label{subsec:finite-typing-separations}


The filtering results of Section~\ref{sec:prelim} give useful sufficient
constructions, but finite typing is not merely regular post-filtering of the
trivial-typing class.  The next linear example gives a stronger separation and,
at the same time, a sharp fan-out effect.

Let
\[
  L_{\mathrm{cd}}:=
  \{\mathtt a^n\mathtt c\mathtt b^n:n\ge0\}
  \cup
  \{\mathtt a^n\mathtt d\mathtt b^{2n}:n\ge0\}
\]
over $\Sigma_{\mathrm{cd}}:=\{\mathtt a,\mathtt b,\mathtt c,\mathtt d\}$.  Let
$M_{\mathrm{cd}}:=\{1,\eta_{\mathtt c},\eta_{\mathtt d},0\}$, where $1$ is the identity, $0$ is a zero,
and every product of two elements of $\{\eta_{\mathtt c},\eta_{\mathtt d},0\}$ is $0$.  Define
\[
 h_{\mathrm{cd}}(\mathtt a)=h_{\mathrm{cd}}(\mathtt b)=1,
 \qquad
 h_{\mathrm{cd}}(\mathtt c)=\eta_{\mathtt c},
 \qquad
 h_{\mathrm{cd}}(\mathtt d)=\eta_{\mathtt d}.
\]

\begin{proposition}[A fan-out-sensitive typed language]
\label{prop:strong-fanout-separation}
The language $L_{\mathrm{cd}}$ is linear and
\[
  L_{\mathrm{cd}}\in\Cmcf{1}{h_{\mathrm{cd}}}.
\]
For every $k,\ell\ge0$,
\[
  L_{\mathrm{cd}}\notin\Cmcf{1}{h_{\Sigma_{\mathrm{cd}}}^{k,\ell}}.
\]
There do not exist a language $L_0\in \YSub{1}$ and a regular language
$L_{\mathrm{reg}}$ such that
$L_{\mathrm{cd}}=L_0\cap L_{\mathrm{reg}}$.  Moreover, for every finite monoid
homomorphism $g:\Sigma_{\mathrm{cd}}^*\to N$,
\[
  L_{\mathrm{cd}}\notin\Cmcf{2}{g}.
\]
\end{proposition}

\begin{proof}
The linear grammar
\[
  S\to C\mid D,\qquad
  C\to\mathtt c\mid\mathtt a C\mathtt b,\qquad
  D\to\mathtt d\mid\mathtt a D\mathtt b\mathtt b
\]
generates $L_{\mathrm{cd}}$.

For $(1,h_{\mathrm{cd}})$-substitutability, let equal-type nonempty factors
$x,y$ share an accepting context.  Type $1$ contains no center symbol, so the
context fixes the branch and side and the count constraint forces $x=y$.
For type $\eta_{\mathtt c}$, writing
$x=\mathtt a^{n_1}\mathtt c\mathtt b^{n_2}$ and
$y=\mathtt a^{n_1'}\mathtt c\mathtt b^{n_2'}$, common acceptance gives
$n_1-n_2=n_1'-n_2'$, the invariant for every other $\mathtt c$-context.
For type $\eta_{\mathtt d}$ the invariant is
$2n_1-n_2=2n_1'-n_2'$, while type $0$ has empty distribution.  Hence
$L_{\mathrm{cd}}\in\Cmcf{1}{h_{\mathrm{cd}}}$.

Fix $k,\ell$ and choose $m\ge\max\{1,k+\ell\}$.  The words
$x=\mathtt a^m\mathtt c\mathtt b^m$ and
$y=\mathtt a^m\mathtt d\mathtt b^{2m}$ have the same
$h_{\Sigma_{\mathrm{cd}}}^{k,\ell}$-type and share the empty context, but
$\mathtt a\blank_1\mathtt b$ accepts $x$ and rejects $y$.  Thus every
prefix--suffix typing fails at fan-out one.

Suppose $L_{\mathrm{cd}}=L_0\cap L_{\mathrm{reg}}$ with
$L_0\in\YSub{1}$ and $L_{\mathrm{reg}}$ regular.  Since
$\mathtt c,\mathtt d\in L_0$, substitutability gives
$\D_{L_0}(\mathtt c)=\D_{L_0}(\mathtt d)$; using
$\mathtt a\mathtt d\mathtt b^2,\mathtt a\mathtt c\mathtt b\in L_0$ then gives
$\D_{L_0}(\mathtt c\mathtt b)=\D_{L_0}(\mathtt c\mathtt b^2)$.
Consequently
\[
 \mathtt a^m\mathtt c\mathtt b^{m+t}\in L_0
 \qquad(m\ge1,\ t\ge0).
\]
For $m$ larger than the state number of a DFA for $L_{\mathrm{reg}}$, pumping
the final $\mathtt b^m$ block of the accepted target word gives some $t>0$
with the same word in $L_{\mathrm{reg}}$, contradicting the assumed
intersection representation.

Finally, for finite $g:\Sigma_{\mathrm{cd}}^*\to N$, choose $1\le i<j$ with
\[
 g(\mathtt a^i)=g(\mathtt a^j),\qquad
 g(\mathtt b^i)=g(\mathtt b^j).
\]
The tuples $(\mathtt a^i,\mathtt b^i)$ and
$(\mathtt a^j,\mathtt b^j)$ have equal componentwise $g$-type and are both
accepted by $\blank_1\mathtt c\blank_2$, whereas
$\blank_1\mathtt d\blank_2\mathtt b^i$ accepts only the first:
$\mathtt a^j\mathtt d\mathtt b^{i+j}\notin L_{\mathrm{cd}}$ for $i<j$.
Thus arity two fails for every finite $g$.
\end{proof}

Hence $L_{\mathrm{cd}}$ separates finite typing from every prefix--suffix
typing and arbitrary regular filtering at fan-out one, while no finite typing
repairs the language once arity two is tested.

\subsection{Local Boundary Control and Arbitrary Regular Filtering}
\label{subsec:comparison-separations}

Following the standard local-testability terminology~\cite{McNaughtonPapert1971}, call
$L_{\mathrm{sl}}\subseteq\Sigma^*$ \emph{strictly 2-local} if there is a set
$\mathsf{Adj}\subseteq(\Sigma\cup\{\vdash,\dashv\})^2$ such that
\[
 w\in L_{\mathrm{sl}}
 \quad\Longleftrightarrow\quad
 \text{every adjacent pair in $\vdash w\dashv$ belongs to $\mathsf{Adj}$},
\]
where $\vdash,\dashv\notin\Sigma$ are fixed boundary markers.

\begin{lemma}[Strictly 2-local constraints are boundary typed]
\label{lem:sl2-boundary-control}
If $L_{\mathrm{sl}}\subseteq\Sigma^*$ is strictly 2-local, then for every $f\ge1$,
\[
  L_{\mathrm{sl}}\in\Cmcf{f}{h_\Sigma^{1,1}}.
\]
\end{lemma}

\begin{proof}
Let equal-$h_\Sigma^{1,1}$-type nonempty tuples $\tuple x,\tuple y$ share an
accepting context $E_0$.  Internal adjacent pairs of each $y_i$ are permitted
because they occur in $E_0[\tuple y]$.  If $E[\tuple x]\in L_{\mathrm{sl}}$,
every adjacent pair of $E[\tuple y]$ is either in fixed material, internal to
a $y_i$, or crosses a component boundary.  These are permitted respectively
by $E[\tuple x]$, by $E_0[\tuple y]$, or by equality of the first/last letters
encoded by the boundary type; end-marker pairs are unchanged.  Thus
$E[\tuple y]\in L_{\mathrm{sl}}$, and symmetry finishes the proof.
\end{proof}

\paragraph{Arbitrary target-dependent regular filters.}
For comparison, on a fixed alphabet define
\[
 \mathcal F_f^{\mathrm{arb}}
 :=\bigl\{L_0\cap R:L_0\in\YSub{f},\ R\text{ regular}\bigr\}
   \cap f\text{-}\mathrm{MCFL}.
\]
This class contains all regular languages
(take $\Sigma^*\in\YSub{f}$) and is therefore
not TxtEx-identifiable as a single
positive-data class by Gold's obstruction.
We use it only as an expressiveness baseline,
not as an additional learnable class.

The intrinsic-rank-two witness lies in this larger
fragment:
\[
 L_{\mathrm{intr}}=\Phi_{\mathrm{intr}}^{-1}(0)\cap Q_{\mathrm{intr}}.
\]
It demonstrates intrinsic rank rather than
separation from arbitrary regular filtering.
Proposition~\ref{prop:strong-fanout-separation}
gives that separation at fan-out one;
the next theorem establishes it for a
finitely typed non-context-free fan-out-two target.

\begin{theorem}[A fan-out-two typed language beyond arbitrary regular filtering]
\label{thm:arbitrary-filter-separation}
Let $\Sigma_\star:=\{\mathtt a,\mathtt b,\mathtt c,\mathtt d\}$ and
\[
 L_\star
 :=\{\mathtt a^n\mathtt c^{n+1}\mathtt b^n:n\ge0\}
   \cup
   \{\mathtt a^n\mathtt d^{n+1}\mathtt b^{2n}:n\ge0\}.
\]
Let
$g_\star:\Sigma_\star^*\to(\mathcal P(\{\mathtt c,\mathtt d\}),\cup)$
be the homomorphism defined by
\[
 g_\star(\mathtt a)=g_\star(\mathtt b)=\varnothing,
 \qquad
 g_\star(\mathtt c)=\{\mathtt c\},
 \qquad
 g_\star(\mathtt d)=\{\mathtt d\},
\]
and put
\[
 h_\star:=(h_{\Sigma_\star}^{1,1},g_\star).
\]
Then $L_\star$ is non-context-free, has minimum fan-out two and minimum
fan-out-two rule rank one, and for every $f\ge2$,
\[
  L_\star\in\Cmcf{f}{h_\star}
  \setminus\mathcal F_f^{\mathrm{arb}}.
\]
Equivalently, there are no $L_0\in\YSub{f}$ and regular $R_{\mathrm{reg}}$ such that
$L_\star=L_0\cap R_{\mathrm{reg}}$.
\end{theorem}

\begin{proof}
A good fan-out-two rank-one grammar is
\[
\begin{aligned}
 A_{\mathtt c}
   &\to(\mathtt a\mathtt c,\mathtt b)
     \mid(\mathtt a\xi_1^1\mathtt c,\xi_1^2\mathtt b)(A_{\mathtt c}),\\
 A_{\mathtt d}
   &\to(\mathtt a\mathtt d,\mathtt b\mathtt b)
     \mid(\mathtt a\xi_1^1\mathtt d,\xi_1^2\mathtt b\mathtt b)(A_{\mathtt d}),\\
 S&\to\mathtt c\mid\mathtt d
   \mid(\xi_1^1\mathtt c\xi_1^2)(A_{\mathtt c})
   \mid(\xi_1^1\mathtt d\xi_1^2)(A_{\mathtt d}).
\end{aligned}
\]
It generates exactly $L_\star$.  The intersection with the regular language
$\mathtt a^*\mathtt c^+\mathtt b^*$ is
$\{\mathtt a^n\mathtt c^{n+1}\mathtt b^n:n\ge0\}$, which is not
context-free.  For a short closure proof, introduce a fresh symbol
$\mathtt\#$ and the homomorphism fixing $\mathtt a,\mathtt b,\mathtt c$ and
sending $\mathtt\#$ to $\mathtt c$.  The inverse image of this branch,
intersected with
$\mathtt a^*\mathtt c^*\mathtt\#\mathtt b^*$, is
$\{\mathtt a^n\mathtt c^n\mathtt\#\mathtt b^n:n\ge0\}$; erasing
$\mathtt\#$ would make $\{\mathtt a^n\mathtt c^n\mathtt b^n:n\ge0\}$
context-free, a contradiction.  Thus the minimum fan-out is two, and
infinitude together with the displayed grammar gives minimum fan-out-two rule
rank one.

For typed substitutability put
\[
 Q_\star:=\mathtt a^*\mathtt c^+\mathtt b^*
          \cup\mathtt a^*\mathtt d^+\mathtt b^*,
 \qquad
 \Psi(w):=(|w|_{\mathtt a},|w|_{\mathtt b},
            |w|_{\mathtt c},|w|_{\mathtt d}).
\]
The strictly $2$-local language $Q_\star$ is boundary typed by
Lemma~\ref{lem:sl2-boundary-control}.  Let equal-$h_\star$-type tuples
$\tuple x,\tuple y$ share an accepting context $E_0$, and set
\[
 \Delta:=\sum_i\Psi(y_i)-\sum_i\Psi(x_i).
\]
Their equal $g_\star$-types put $E_0[\tuple x],E_0[\tuple y]$ on the same
branch, so for some $q\in\mathbb Z$,
\[
 \Delta=q(1,1,1,0)\quad\text{or}\quad\Delta=q(1,2,0,1)
\]
according as that branch is $\mathtt c$ or $\mathtt d$.

If $F[\tuple x]\in L_\star$, boundary typing gives
$F[\tuple y]\in Q_\star$, and equal $g_\star$-types put the two fillings on
the same $F$-branch.  When $E_0$ and $F$ choose the same branch, the displayed
count difference preserves its affine relation.  If they choose different
branches, no tuple component can contain $\mathtt c$ or $\mathtt d$; hence the
two center coordinates of $\Delta$ vanish, forcing $q=0$.  In either case
$F[\tuple y]\in L_\star$, and symmetry proves
$L_\star\in\Cmcf{f}{h_\star}$.

Finally assume $L_\star=L_0\cap R_{\mathrm{reg}}$ with
$L_0\in\YSub{f}$ and $R_{\mathrm{reg}}$ regular.  Since
$\mathtt c,\mathtt d\in L_0$, arity-one substitutability gives
$\D_{L_0}(\mathtt c)=\D_{L_0}(\mathtt d)$.  Hence, replacing the two occurrences of $\mathtt d$ successively,
$\mathtt a\mathtt d^2\mathtt b^2\in L_0$ yields first
$\mathtt a\mathtt c\mathtt d\mathtt b^2\in L_0$ and then
$\mathtt a\mathtt c^2\mathtt b^2\in L_0$; together with
$\mathtt a\mathtt c^2\mathtt b\in L_0$ gives
\[
 \D_{L_0}(\mathtt c^2\mathtt b)
 =\D_{L_0}(\mathtt c^2\mathtt b^2).
\]
Thus $\mathtt a^m\mathtt c^{m+1}\mathtt b^{m+t}\in L_0$ for all
$m\ge1,t\ge0$.  For $m$ at least the state number of a DFA for
$R_{\mathrm{reg}}$, pumping the final $\mathtt b^m$ block of the accepted
target word gives some $t>0$ with the same word in $R_{\mathrm{reg}}$, hence
in the intersection but not in $L_\star$, a contradiction.
\end{proof}

Thus the single target $L_\star$ rules out representation by a Yoshinaka
language followed by any finite-state filter, irrespective of how the filter
typing refines $h_\star$.

This rank-one separation complements
Theorem~\ref{thm:intrinsic-ranktwo-witness},
which gives a boundary-typed non-context-free
target of minimum rule rank two across all finite fan-outs.

\section{Conclusion}
\label{sec:conclusion}

Finite-monoid typing gives a compositional bias for positive-data MCFG
learning: every fixed $(f,h)$ admits exact reconstruction and TxtEx
identification, with polynomial-time updates under a fixed rule-rank cap.
Fan-out is part of the bias: $L_{\mathrm{cd}}$
shows that increasing tested arity can destroy
substitutability beyond repair by any finite typing.
The language $L_\star$ additionally excludes
arbitrary target-dependent regular filtering.

For characteristic data, arbitrary fixed typing has a grammar-size polynomial
at rule rank one, while prefix--suffix-determined typing has a
grammar-size-and-rule-thickness polynomial at arbitrary fixed rank.
Theorem~\ref{thm:intrinsic-ranktwo-witness} supplies a boundary-typed
non-context-free target of minimum fan-out and rule rank two.  Full yield-type
refinement and empty-component elimination can nevertheless cause exponential
thickness growth; Appendix~\ref{app:normalization-details} records the latter
normalization details.

\begin{openproblem}[Original-thickness bounds beyond fan-out one]
\label{op:fanout-two}
Fix $f\ge2$ and a rule-rank bound $r\ge2$.
Does there exist a fixed finite typing $h$ and a family of
reduced good fan-out-$\le f$, rank-$\le r$ grammars $G_n$
of targets in $\Cmcf{f}{h}$ for which
$\|G_n\|\max\{1,\theta_{G_n}\}$ grows polynomially
in $n$ but every characteristic sample for
$\mathcal B^{\le r}_{f,h}$ has superpolynomial
total positive length?  The fan-out-one counterexample of the companion
CFG paper~\cite{kuriyamaCFG} does not by itself
answer this higher-arity question:
membership in $\Cmcf{1}{h}$ need not imply
membership in $\Cmcf{f}{h}$ for $f\ge2$.
\end{openproblem}



The window--abelian-group extension is proved separately in
Appendix~\ref{app:window-abelian}. The analogous
higher-fan-out result for nonabelian groups
is not established here.

Other directions include adaptive learning over several fan-out biases,
learning useful typings without an a priori codomain bound, characterizing the
rank-one subfamily inside the typed classes, understanding the relation with
full-multicontext congruential MCFGs, and extending the reconstruction to
PMCFG-style copying.

\section*{Acknowledgments}
The author thanks Makoto Kanazawa and Ryo Yoshinaka for valuable discussions
and encouragement.

\section*{Declaration of competing interest}
The author declares that he has no known competing financial interests or
personal relationships that could have appeared to influence the work reported
in this paper.

\section*{Funding}
This research received no specific grant from any funding agency in the public,
commercial, or not-for-profit sectors.

\section*{Data availability}
No data were used for the research described in this article.

\section*{Declaration of generative AI and AI-assisted technologies in the manuscript preparation process}
The author used ChatGPT (OpenAI) and Claude (Anthropic) during the research
and manuscript-preparation process, including exploratory mathematical
argument development, proof drafting and criticism, literature and
consistency checks, and editing of text and LaTeX source.
AI-generated suggestions and reviews are not treated as independent
mathematical verification. The author takes responsibility for the
mathematical claims, citations, and final content of this article.

\appendix
\section{Details of the Good Normalization}
\label{app:normalization-details}

The following bookkeeping records the fixed-parameter bounds used in
Proposition~\ref{prop:good-normalization}; semantic correctness of the
nondeleting and nonmerging conversions is standard
\cite{SekiEtAl1991,Kanazawa2019Ogden,Yoshinaka2011}.

\paragraph{Nondeleting conversion and component permutations.}
After trimming, use only the nondeleting step (f3) from
the proof of Seki et al.'s Lemma~2.2, in Kanazawa's rank-preserving form
\cite[proof of Lemma~2.2, item~(f3)]{SekiEtAl1991}\cite[Section~2.1]{Kanazawa2019Ogden}.
Retained-component subsets give at most $O_f(1)$ decorated copies per
nonterminal and $O_{f,r_{\max}}(1)$ projected copies per rule; deleting an
empty child cannot increase rank.  For nonpermutation, decorate a fan-out-$d$
nonterminal by $\sigma\in\mathfrak S_d$.  The parent permutation determines
the child permutations, so there are at most $f!$ copies, no child is
duplicated, and rank is preserved.  Denote the result by $G_{\mathrm{nd}}$.

\paragraph{Removing empty tuple components.}
For a fan-out-$d$ nonterminal $A$, let
\[
 \mathsf{Pat}(A):=
 \{I\subseteq\{1,\ldots,d\}:\text{$I$ is the nonempty-component pattern of
 some }\tuple w\in L_A(G_{\mathrm{nd}})\}.
\]
Finite saturation computes these sets.  For realizable child patterns, replace
variables from absent components by $\eps$, delete empty parent components,
and omit all-empty children.  There are at most $2^f$ patterns per nonterminal
and $2^{fr_{\max}}$ choices per rule.  Introduce $A^I$ for each nonempty
pattern $I$.  Retained variables still occur once and in order, so
nondeletion/nonpermutation, fan-out, and rank are preserved and the result is
$\lambda$-free.

Induction on derivations gives
$L_{A^I}$ exactly as the projection of source tuples having pattern $I$;
all-empty omitted children can be restored by terminal derivations realizing their all-empty component pattern.  Since
the start symbol has fan-out one, $S^{\{1\}}$ generates precisely the nonempty
part of the original language.  Projection shortens templates, giving
$O_{f,r_{\max}}(\|G_{\mathrm{nd}}\|)$ size.

\paragraph{Nonmerging.}
Apply Yoshinaka's nonmerging conversion
\cite[Section~2.3, Lemma~1]{Yoshinaka2011}.  Each decorated nonterminal denotes exactly the projected tuples associated
with its retained component pattern, and adjacent-component concatenation
preserves $\lambda$-freeness, nonpermutation, fan-out, and rule rank.
For fixed fan-out and rank Yoshinaka's combined good-form conversion has
$O(\|G\|)$ size~\cite[Section~2.3, after Example~2]{Yoshinaka2011}; trimming
can only decrease all bounds.  This proves
Proposition~\ref{prop:good-normalization}.

\subsection*{Thickness transfer for $\lambda$-free presentations}

\begin{proposition}[Thickness under good-form conversion]
\label{prop:structural-good-thickness}
Fix $f$ and $r_{\max}$.  Let $G_{\mathrm{in}}$ be a reduced $\lambda$-free
$f$-MCFG of maximum rule rank at most $r_{\max}$, possibly with deleting
rules, and suppose $L(G_{\mathrm{in}})\cap\Sigma^+\neq\emptyset$.  The good
normalization of Proposition~\ref{prop:good-normalization} yields an equivalent
reduced good grammar $G$ with the same fan-out and rule-rank bounds such that
\[
  \|G\|=O_{f,r_{\max}}(\|G_{\mathrm{in}}\|),\qquad
  t_G\le t_{G_{\mathrm{in}}},\qquad
  \theta_G=O_{f,r_{\max}}\!\bigl(
     \|G_{\mathrm{in}}\|\max\{1,t_{G_{\mathrm{in}}}\}\bigr).
\]
\end{proposition}

\begin{proof}
For $\lambda$-free $G_{\mathrm{in}}$, every derived component is nonempty.
The nondeleting conversion therefore needs no later empty-pattern split, and
each retained decoration $A^I$ generates exactly the projection of
$L_A(G_{\mathrm{in}})$ onto $I$.  Hence its minimum total yield cannot exceed
that of $A$.  Component permutations preserve total length, and Yoshinaka's
nonmerging copies are exact adjacent-component images
\cite[Section~2.3, Lemma~1]{Yoshinaka2011}, so
$t_G\le t_{G_{\mathrm{in}}}$.  All decorations have fixed-parameter
linear overhead, giving
$\|G\|=O_{f,r_{\max}}(\|G_{\mathrm{in}}\|)$; then
Lemma~\ref{lem:thickness-equivalence} yields the displayed bound on
$\theta_G$.
\end{proof}

Thus deleting rules do not by themselves obstruct polynomial thickness
transfer from a reduced $\lambda$-free presentation; genuine empty components
are the problematic step.

\begin{proposition}[Empty-component elimination can blow up thickness]
\label{prop:empty-component-blowup}
For every $n\ge1$ there is a reduced fan-out-two, rule-rank-two MCFG $E_n$ of
size $O(n)$ with $t_{E_n}\le2$ and $\theta_{E_n}=O(n)$ such that the
empty-component elimination stage of Proposition~\ref{prop:good-normalization}
contains a retained nonterminal and a retained rule of thickness at least
$2^n$.
\end{proposition}

\begin{proof}
Let $E_n$ have fan-out-one $S$, fan-out-two $A_0,\ldots,A_n$, and
\[
\begin{aligned}
 A_0&\to(\mathtt a,\eps),\\
 A_i&\to(\eps,\mathtt c)\mid
      (\xi_1^1\xi_2^1,\xi_1^2\xi_2^2)(A_{i-1},A_{i-1})
      \quad(1\le i\le n),\\
 S&\to(\xi_1^1\mathtt d\,\xi_1^2)(A_n).
\end{aligned}
\]
It is reduced, nondeleting, nonpermuting, and nonmerging.  Short side
derivations give $t_{E_n}\le2$, and placing any chosen rule on one branch
of the binary spine while completing siblings minimally gives
$\theta_{E_n}=O(n)$.

If an $A_i$-tuple has nonempty first component and empty second component, the
rule $(\eps,\mathtt c)$ cannot occur.  Inductively the shortest such tuple is
$(\mathtt a^{2^i},\eps)$.  Thus pattern $\{1\}$ for $A_n$ survives
empty-component elimination, and the projected start rule
\[
 S\to\xi_1^1\mathtt d\,(A_n^{\{1\}})
\]
can succeed only with a child of length at least $2^n$.  Hence the retained
nonterminal and rule have thickness at least $2^n$.
\end{proof}

This is the same obstruction exhibited by
Proposition~\ref{prop:typed-thickness-exponential}: finite state refinement can
force exponentially long shortest witnesses inside one refined state.

\section{Relation to Yoshinaka's General Framework}
\label{sec:general-framework}

The bounded-rank learner is a typed instance of
Yoshinaka's substitutable-signature framework
\cite{Yoshinaka2015General}. We make the
correspondence explicit to distinguish prior
qualitative learning results from the total-length
bounds of Section~\ref{sec:char-data}.

Fix finite $\Sigma$, $f,r\ge1$ and
$h:\Sigma^*\to M$, and put $M^+:=h(\Sigma^+)$.
The signature $\mathfrak S_{f,r,h}$ has a
start sort $O_*=\Sigma^*$ and tagged sorts
\[
 T_{d,\boldsymbol\mu}
 =\{\tuple w\in(\Sigma^+)^d:
 h^{(d)}(\tuple w)=\boldsymbol\mu\},
 \qquad
 1\le d\le f,\quad\boldsymbol\mu\in(M^+)^d.
\]
Its function names comprise all well-typed
good MCFG templates of child rank at most $r$
(including rank zero); promotions
$T_{1,\mu}\to O_*$; and a nullary $\eps$
constant of sort $O_*$. No function takes an
$O_*$ input. After forgetting the type tags, this is a
finite sort refinement, in the sense of
\cite[Section~8]{Yoshinaka2015General}, of
the signature of good $\lambda$-free MCFG
operations on nonempty tuples. That
untyped signature is a nonempty-component
variant of \cite[Example~3]{Yoshinaka2015General},
rather than literally the signature of that
example; at fan-out one the prefix--suffix
case corresponds to
\cite[Example~4]{Yoshinaka2015General}.
Write $\mathcal S_O,\mathcal C_O$ for the
substructures and contexts realizable in
some grammar over this signature.

\begin{lemma}[Signature contexts are admissible sentence contexts]
\label{lem:signature-contexts}
For every non-start sort,
\[
 \mathcal S_{T_{d,\boldsymbol\mu}}
 =T_{d,\boldsymbol\mu},
 \qquad
 \mathcal C_{T_{d,\boldsymbol\mu}}
 =\{\tuple x\mapsto E[\tuple x]:E\in\mathcal E_d\}.
\]
The latter equality is between \emph{functions}
on $T_{d,\boldsymbol\mu}$.
Also $\mathcal C_{O_*}$ contains only the identity.
\end{lemma}

\begin{proof}
Every element of a non-start sort is a rank-zero
good template constant. No operation takes
$O_*$ as input. For a context of non-start
sort, induct along the path from its hole
to the start. Composing a good template with
an admissible parent context preserves the
order and single occurrence of the
distinguished child's variables.
Consecutive holes in one output component
are separated by a terminal or a nonempty
sibling component (nonmerging); holes in
different output components are separated
by a nonempty parent-context separator.
Thus every resulting sentence context
belongs to $\mathcal E_d$, by repeated
application of Lemma~\ref{lem:filling-identity}.

Conversely, for any
$E=v_0\blank_1v_1\cdots\blank_dv_d
\in\mathcal E_d$, the unary good template
$(v_0\xi_1^1v_1\cdots\xi_1^dv_d)$
with a rank-zero child and start promotion
realizes $\tuple x\mapsto E[\tuple x]$.
Here the internal $v_i$ are nonempty,
as required by goodness; the output type
is fixed by $h$.
\end{proof}

\begin{proposition}[Bounded-rank coverage by Yoshinaka's 2015 framework]
\label{prop:yoshinaka2015-embedding}
The signature $\mathfrak S_{f,r,h}$ satisfies
Assumptions~1--3 of
\cite{Yoshinaka2015General}, and
\[
 \mathsf{Subst}(\mathfrak S_{f,r,h})
 =\Cmcf{f}{h}\cap\mathsf{MCFL}(f,r).
\]
Consequently Yoshinaka's identification theorem
applies to the bounded-rank class
\cite[Theorem~1]{Yoshinaka2015General}:
for fixed $f,r,h$, updates are polynomial
in the accumulated sample, and characteristic
sets have cardinality at most
$|N^\sharp|+|P^\sharp|$ for a trimmed
signature presentation
\cite[Lemma~3]{Yoshinaka2015General}.
This does not bound the total \emph{length}
of those samples.
\end{proposition}

\begin{proof}
Every signature grammar is a fan-out-$\le f$,
rank-$\le r$ MCFG after forgetting tags.
Conversely, rank-preserving good normalization
(Proposition~\ref{prop:good-normalization}),
trimmed yield-type refinement
(Definition~\ref{def:output-refinement}),
start promotions and the separate $\eps$
constant present every language in
$\mathsf{MCFL}(f,r)$ over this signature.
At the start sort substitutability is
vacuous; at all other sorts
Lemma~\ref{lem:signature-contexts} identifies
the four-instance signature implication
with Definition~\ref{def:tuple-substitutable}
for each arity and componentwise type.
The trimmed presentation has no useless
nonterminals or functions, as required in
\cite[Section~6]{Yoshinaka2015General}.
This proves the equality.

For the effective assumptions, sort/context
membership and template evaluation are
polynomial-time by the lemma and finite
typing; uniform bounded-rank MCFG membership
is polynomial
\cite[Proposition~4]{Yoshinaka2011}.
All operations have arity at most $r$.
For a sentence of length $n$, enumerate
the start word, each tagged tuple of up
to $f$ nonoverlapping nonempty intervals
with nonempty internal gaps, and the
associated context functions. For each
parent tuple occurrence, enumerate the
rank-$\le r$ good templates determined by
ordered child-component intervals (at most
$fr$ such intervals) and the remaining
terminal material, recording the actual
\emph{function names}. Add promotions and
the $\eps$ constant as appropriate.
By Lemma~\ref{lem:composition-witness-materialization},
every realizable function instance occurs.
At most $2f(r+1)$ cut positions are needed,
so the enumeration is polynomial in $n$
for fixed $f,r,h$. These are precisely
Assumptions~1--3; applying Yoshinaka's
Theorem~1 and Lemma~3 gives the stated
learning consequences.
\end{proof}

\begin{lemma}[Substructure-indexed and guarded reconstruction]
\label{lem:substp}
For the preceding signature and enumerations,
Yoshinaka's $\mathrm{SUBST}_{\mathrm P}(K)$
and $\mathcal B^{\le r}_{f,h}(K)$ generate
the same language for every finite
positive sample $K\subseteq\Sigma^*$.
\end{lemma}

\begin{proof}
Identify their non-start symbols indexed
by the same observed tuples and typed sorts.
A $\mathrm{SUBST}_{\mathrm P}$ rule
with parent $s_0$ and template value
$z=\rho(s_1,\ldots,s_t)$ requires a
common $c\in\mathrm{CON}_O(K)$ such that
$c\odot s_0,c\odot z\in K$,
in the notation of
\cite{Yoshinaka2015General}. By Lemma~\ref{lem:signature-contexts}
this is an admissible sentence context.
The guarded Rule~(Link) sends $[s_0]$
to $[z]$, followed by Rule~(Comp), or
Rule~(Const) for $t=0$, whose witness
is materialized by
Lemma~\ref{lem:composition-witness-materialization}.
Conversely, every guarded Link is a
$\mathrm{SUBST}_{\mathrm P}$ instance of
the unary identity template; every Comp
and Const is its own witnessed parent
instance. Induction on derivations
therefore equates the non-start
symbol languages.

At the start sort the only context is
the identity and no operation accepts
a start-sort argument. Each observed
start symbol consequently has the
same language
\[
 \bigcup_{w\in K\cap\Sigma^+}L([(w)])
 \ \cup\ (\{\eps\}\cap K),
\]
by the promotions and direct $\eps$
constant. This is precisely guarded
Rule~(Start); for $K=\emptyset$ both
languages are empty.
\end{proof}

\paragraph{Characteristic witness correspondence.}
Yoshinaka's proof uses the context-indexed
$\mathrm{SUBST}_{\mathrm D}$, whereas its
Section~7 discusses $\mathrm{SUBST}_{\mathrm P}$
under the stated additional assumption
\cite{Yoshinaka2015General}.
The sample-level equality here is established
directly by Lemma~\ref{lem:substp}.
Choosing its non-start witnesses as
$c_X=\chi(X)$ and $s_X=\omega(X)$
gives precisely $\Wit(\widetilde G_0)$
at non-start sorts; promoted start
witnesses are among these chosen words,
and $\eps$ contributes its direct
datum when present. This identifies
the positive-word witness \emph{pattern}
with consistent start conventions,
not the syntactic hypothesis grammars.

% #2 MCFG v72 日本語査読用付録
% 数式・定理 label・証明の依存関係は英語版 appendix_window_abelian.tex と対応させる．
\section{固定窓と有限アーベル群を組み合わせた型付け}
\label{app:window-abelian}

この付録では，固定接頭辞--接尾辞窓と有限アーベル群の情報を組み合わせた場合に，
特徴データの総文字数についての多項式定理を拡張する．
群を用いない元の定理~\ref{thm:boundary-polydata} は変更しない．

\begin{lemma}[アーベル群値を指定した短いタプル]
\label{lem:group-typed-shortest}
$G=(V,\Sigma,P,S,\fan)$ を，ファンアウトが高々 $f$，
規則ランクが高々 $r_{\max}\ge1$ の縮約済み良形 MCFG とする．
$g:\Sigma^*\to A$ を，固定有限アーベル群への準同型とする．
各 $B\in V$ と，実現可能な任意のベクトル
$\mathbf s\in A^{\fan(B)}$ に対し，
\[
 \tuple u\in L_B(G),\qquad
 g^{(\fan(B))}(\tuple u)=\mathbf s
\]
を満たすタプルで，
\[
 \|\tuple u\|_1\le\beta_g:=
 (r_{\max}+1)\bar\theta_G
 \max\{1,(|A|^f-1)|V||A|^f\}
\]
となるものが存在する．
\end{lemma}

\begin{proof}
$B$ からの終端導出で，根の $g$-値ベクトルが $\mathbf s$ となるものを考える．
その中から，まず根の生成タプルの総文字数を最小化し，
次に導出木の節点数を最小化した導出 $T$ を選ぶ．
各節点に，その非終端記号 $C$ と，その部分木が生成するタプルの
成分ごとの $g$-値ベクトルを標識する．
この標識は高々 $n_g:=|V||A|^f$ 種類しかない．

ある祖先とその真の子孫が同じ標識をもつと仮定する．
祖先の部分木を子孫の部分木で置き換えると，
根の $g$-値ベクトルが保存される．
実際，MCFG の各規則は，子の成分別群値から親の群値を計算する．
非消去性により生成タプルの総文字数は増えず，
導出木の節点数は真に減る．
これは $T$ の選び方に矛盾する．
したがって木の深さは高々 $n_g$ である．

部分木が生成するタプルの総文字数が
$\bar\theta_G$ を超える節点を\emph{重い節点}と呼ぶ．
非消去性により，重い節点の祖先もすべて重い．
互いに比較不能な重い節点を
$v_1,\ldots,v_m$ とし，根のファンアウトを $d\le f$ とする．
各 $v_i$ の標識を $C_i$ とする．
補題~\ref{lem:thickness-local-bounds}(ii) により，
$C_i$ から長さ高々 $\bar\theta_G$ のタプルを生成する導出が存在する．
$v_i$ の部分木をこの短い導出で置き換えれば，
根の総文字数は真に減る．

文法は線形かつ非消去であるから，互いに比較不能な部分木の出力成分は，
根の各成分内で互いに交わらない終端ブロックに寄与する．
ただし成分ごとにその出現順序が異なることはある．
アーベル群では積を加法的に表せ，ブロックの順序に依存しない．
したがって，各節点の置換による変化を
$\delta_i\in A^d$ とおくと，
任意の部分集合 $J$ に対応する同時置換は
根の $g$-値ベクトルを $\sum_{i\in J}\delta_i$ だけ変化させる．
$m\ge|A|^d$ なら，$m+1$ 個の部分和
\[
 0,\quad\delta_1,\quad
 \delta_1+\delta_2,\quad\ldots,\quad
 \delta_1+\cdots+\delta_m
\]
のうち二つが一致する．
よって非空の連続部分集合に対応する置換は根の群値を保つ一方，
生成タプルを真に短くする．これは最小性に矛盾する．
したがって重い節点からなる任意の反鎖の大きさは
$|A|^d\le|A|^f$ 未満である．

根が軽ければ生成タプルの総文字数は
$\bar\theta_G$ 以下である．
根が重い場合，重い節点たちは祖先について閉じた木をなす．
その葉は重い節点の反鎖であり，各重い節点はこのいずれかの葉への経路上にある．
各経路の長さは高々 $n_g$ なので，
重い節点は高々 $(|A|^f-1)n_g$ 個である．
補題~\ref{lem:thickness-local-bounds}(i) により，
各規則が明示的に書く終端因子の総長は
$\bar\theta_G$ 以下である．
重い節点の軽い子はそれぞれ高々
$\bar\theta_G$ 個の終端記号に寄与し，
重い節点の子の個数は高々 $r_{\max}$ である．
重い節点とその軽い子を数えれば所望の上界が得られる．
$|A|=1$ の場合は，根が軽い場合の議論で扱える．
\end{proof}

\begin{lemma}[窓・アーベル群型付けの短い正準タプル]
\label{lem:window-abelian-short}
補題~\ref{lem:group-typed-shortest} の仮定の下で，
$k,\ell\ge0$ を固定する．
$h:\Sigma^*\to M$ を
$h\preceq(h_\Sigma^{k,\ell},g)$
を満たす有限型付けとする．
$\widetilde G_0$ を $G$ の完全 $h$-型付き精密化を刈り込んだものとし，
その開始記号以外の非終端記号数を $n_{\mathrm{typ}}$ とする．
このとき各非終端記号 $X$ について，正準タプル
$\omega(X)$ を
\[
 \|\omega(X)\|_1\le
 f(k+\ell+1)n_{\mathrm{typ}}\,
 \bigl(\bar\theta_G+r_{\max}\beta_g\bigr)
\]
となるように選べる．
\end{lemma}

\begin{proof}
$X$ からの型付き終端導出のうち，節点数が最小のものを選ぶ．
一つの経路上で型付き非終端記号が反復すれば，
根の $h$-型を保ったまま反復部分を除けるので，
各経路の節点数は高々 $n_{\mathrm{typ}}$ である．
補題~\ref{lem:boundary-short-canonical} の証明と同様に，
根の各成分で $h_\Sigma^{k,\ell}$ の型を決める位置を標識する．
成分が窓の閾値より短い場合は全位置を標識し，
それ以外は先頭 $k$ 位置と末尾 $\ell$ 位置を標識する．
$k=\ell=0$ の場合は1位置を標識する．
これらの標識への祖先経路がつくる骨格の節点数は高々
$f(k+\ell+1)n_{\mathrm{typ}}$ である．

骨格に含めなかった子部分木は，
それが寄与する根の各成分において，
標識されていない内部部分に完全に収まる．
その非終端記号 $C$ と，実際の成分別群値
$\mathbf s$ を記録する．
補題~\ref{lem:group-typed-shortest} により，
$C$ から同じ $\mathbf s$ を生成し，
総文字数が $\beta_g$ 以下の導出で置き換える．
良形性により置換後の全成分は非空に保たれる．
保護された境界位置は変化せず，
もとの長い成分は置換後も長いままである．
よって窓型は保存される．
また，根の各成分は規則中の明示的な終端因子と子成分の連接からなる．
置換した子成分の群値はそれぞれ保存されているので，
根の $g$-値ベクトルも変化しない．
この段階では群の可換性を新たに使う必要はない．

したがって積型 $(h_\Sigma^{k,\ell},g)$ の値が保存され，
それを通じて決定されるより粗い $h$-型も保存される．
中間の型付き標識はいったん忘れ，
得られた $G$ 上の導出を
命題~\ref{prop:output-invariants}(ii) で
$h$-型付き精密化に一意に持ち上げる．
$X$ は刈り込み後に残る非終端記号であるため，
この根をもつ導出も刈り込み後の文法に残る．
骨格の各節点が明示的に書く終端記号は
高々 $\bar\theta_G$ 個であり，
骨格から外れた子は高々 $r_{\max}$ 個である．
これらを数えて表示した長さ上界を得る．
\end{proof}

\begin{theorem}[窓・アーベル群型付けに対する多項式特徴データ]
\label{thm:window-abelian-polydata}
$f,r_{\max},k,\ell,\Sigma$，
有限アーベル群 $A$ と準同型 $g:\Sigma^*\to A$，
および有限型付け $h\preceq(h_\Sigma^{k,\ell},g)$ を固定する．
ファンアウトが高々 $f$，規則ランクが高々 $r_{\max}$ の
縮約済み良形文法 $G$ が
$\Cmcf{f}{h}$ に属する言語を生成するとする．
このとき，ランクを制限した再構成作用素
$\mathcal B^{\le r_{\max}}_{f,h}$ は，
次を満たす特徴標本をもつ．
\[
 |\Wit(\widetilde G_0)|=O(\|G\|),\qquad
 \|\Wit(\widetilde G_0)\|_+
   =O_{f,r_{\max},k,\ell,h,g,\Sigma}
     (\|G\|^3\bar\theta_G).
\]
\end{theorem}

\begin{proof}
固定有限型付けにより
$n_{\mathrm{typ}}\le|V||M|^f=O(\|G\|)$ である．
補題~\ref{lem:group-typed-shortest} から
$\beta_g=O(\|G\|\bar\theta_G)$ が従う．
補題~\ref{lem:window-abelian-short} により，
選択した正準タプルの総文字数は
$O(\|G\|^2\bar\theta_G)$ 以下である．
補題~\ref{lem:untyped-contexts} を適用すると，
各特徴証拠語の長さも
$O(\|G\|^2\bar\theta_G)$ 以下となる．
命題~\ref{prop:witness-cardinality} から
その語数は $O(\|G\|)$ である．
空語 $\eps$ の証拠が必要でも，
符号化された総長の増加は高々定数である．
\end{proof}

\begin{remark}
可換性を使うのは，
補題~\ref{lem:group-typed-shortest} の反鎖に対する零和の段階である．
ファンアウト1では各部分木は一つの区間に寄与し，
関連する CFG 論文は任意の有限群を扱っている．
一方，$f\ge2$ で群が非可換な場合，
異なる部分木の寄与が根の各成分で異なる順序に交錯し得る．
本稿の証明はこの場合を扱わないが，
不可能性を主張するものではない．
\end{remark}


\small
\begin{thebibliography}{99}

\bibitem{Gold1967}
E.~M.~Gold.
\newblock Language identification in the limit.
\newblock {\em Information and Control}, 10(5):447--474, 1967.

\bibitem{Angluin1980}
D.~Angluin.
\newblock Inductive inference of formal languages from positive data.
\newblock {\em Information and Control}, 45(2):117--135, 1980.

\bibitem{ClarkEyraud2007}
A.~Clark and R.~Eyraud.
\newblock Polynomial identification in the limit of substitutable context-free
languages.
\newblock {\em Journal of Machine Learning Research}, 8:1725--1745, 2007.

\bibitem{deLaHiguera1997}
C.~de~la~Higuera.
\newblock Characteristic sets for polynomial grammatical inference.
\newblock {\em Machine Learning}, 27:125--138, 1997.

\bibitem{CosteTyping2004}
F.~Coste, D.~Fredouille, C.~Kermorvant, and C.~de~la~Higuera.
\newblock Introducing domain and typing bias in automata inference.
\newblock In \emph{Grammatical Inference: Algorithms and Applications -- ICGI 2004},
LNCS 3264, pp.~115--126, Springer, 2004.

\bibitem{Takada1988}
Y.~Takada.
\newblock Grammatical inference for even linear languages based on control sets.
\newblock \emph{Information Processing Letters}, 28(4):193--199, 1988.
\newblock doi:10.1016/0020-0190(88)90208-6.

\bibitem{Takada1995}
Y.~Takada.
\newblock Learning formal languages based on control sets.
\newblock In K.~P. Jantke and S.~Lange (eds.), \emph{Algorithmic Learning
for Knowledge-Based Systems}, LNCS 961, pp.~316--339, Springer, 1995.

\bibitem{Clark2013}
A.~Clark.
\newblock Learning trees from strings: A strong learning algorithm for some context-free grammars.
\newblock \emph{Journal of Machine Learning Research}, 14:3537--3559, 2013.

\bibitem{CosteGaretNicolas2012}
F.~Coste, G.~Garet, and J.~Nicolas.
\newblock Local substitutability for sequence generalization.
\newblock In {\em Proceedings of the 11th International Conference on
Grammatical Inference}, JMLR Workshop and Conference Proceedings 21,
pp.~97--111, 2012.

\bibitem{EyraudHeinzYoshinaka2016}
R.~Eyraud, J.~Heinz, and R.~Yoshinaka.
\newblock Efficiency in the identification in the limit learning paradigm.
\newblock In J.~Heinz and J.~M.~Sempere (eds.), {\em Topics in Grammatical
Inference}, pp.~25--46. Springer, 2016.
\newblock doi:10.1007/978-3-662-48395-4\_2.

\bibitem{CosteNicolas2019}
F.~Coste and J.~Nicolas.
\newblock Learning local substitutable context-free languages from positive
examples in polynomial time and data by reduction.
\newblock In {\em Proceedings of the 14th International Conference on
Grammatical Inference}, Proceedings of Machine Learning Research 93,
pp.~155--168, 2019.

\bibitem{Yoshinaka2008}
R.~Yoshinaka.
\newblock Identification in the limit of $k,l$-substitutable context-free
languages.
\newblock In {\em Grammatical Inference: Algorithms and Applications},
LNCS 5278, pp.~266--279. Springer, 2008.

\bibitem{Yoshinaka2009}
R.~Yoshinaka.
\newblock Learning mildly context-sensitive languages with multidimensional
 substitutability from positive data.
\newblock In {\em Algorithmic Learning Theory}, LNCS 5809, pp.~278--292.
 Springer, 2009.
\newblock doi:10.1007/978-3-642-04414-4\_24.

\bibitem{Yoshinaka2011}
R.~Yoshinaka.
\newblock Efficient learning of multiple context-free languages with
multidimensional substitutability from positive data.
\newblock {\em Theoretical Computer Science}, 412(19):1821--1831, 2011.
\newblock doi:10.1016/j.tcs.2010.12.058.

\bibitem{Yoshinaka2015General}
R.~Yoshinaka.
\newblock General perspective on distributionally learnable classes.
\newblock In \emph{Proceedings of the 14th Meeting on the Mathematics
of Language (MoL 2015)}, pp.~87--98.
Association for Computational Linguistics, 2015.
\newblock doi:10.3115/v1/W15-2308.

\bibitem{Yoshinaka2010}
R.~Yoshinaka.
\newblock Polynomial-time identification of multiple context-free languages
from positive data and membership queries.
\newblock In {\em Grammatical Inference: Theoretical Results and Applications},
LNCS 6339, pp.~230--244. Springer, 2010.
\newblock doi:10.1007/978-3-642-15488-1\_19.

\bibitem{ClarkYoshinaka2016}
A.~Clark and R.~Yoshinaka.
\newblock Distributional learning of context-free and multiple context-free grammars.
\newblock In J.~Heinz and J.~M.~Sempere (eds.), \emph{Topics in Grammatical
Inference}, pp.~143--172. Springer, 2016.
\newblock doi:10.1007/978-3-662-48395-4\_6.

\bibitem{ClarkYoshinaka2014}
A.~Clark and R.~Yoshinaka.
\newblock Distributional learning of parallel multiple context-free grammars.
\newblock {\em Machine Learning}, 96(1--2):5--31, 2014.
\newblock doi:10.1007/s10994-013-5403-2.

\bibitem{OatesEtAl2006}
T.~Oates, T.~Armstrong, L.~Becerra-Bonache, and M.~Atamas.
\newblock Inferring grammars for mildly context-sensitive languages in polynomial-time.
\newblock In {\em Grammatical Inference: Algorithms and Applications -- ICGI 2006},
LNCS 4201, pp.~137--147. Springer, 2006.

\bibitem{KanazawaYoshinaka2023}
M.~Kanazawa and R.~Yoshinaka.
\newblock Extending distributional learning from positive data and membership queries.
\newblock In {\em Proceedings of the 16th International Conference on Grammatical Inference},
Proceedings of Machine Learning Research 217, pp.~8--22, 2023.

\bibitem{Gallot2021}
P.~D.~Gallot.
\newblock \emph{Safety of transformations of data trees: Tree transducer
theory applied to a verification problem on shell scripts}.
\newblock PhD thesis, Universit\'e de Lille, 2021.
\newblock
\url{https://pepite-depot.univ-lille.fr/LIBRE/EDMADIS/2021/2021LILUB030.pdf}.

\bibitem{BishopEtAl2026EDT0L}
A.~Bishop, M.~Elder, A.~Evetts, P.~Gallot, and A.~Levine.
\newblock On groups with EDT0L word problem.
\newblock \emph{International Journal of Algebra and Computation},
36(4):425--486, 2026.
\newblock doi:10.1142/S0218196726500165.

\bibitem{RambowSatta1999}
O.~Rambow and G.~Satta.
\newblock Independent parallelism in finite copying parallel rewriting systems.
\newblock {\em Theoretical Computer Science}, 223(1--2):87--120, 1999.
\newblock doi:10.1016/S0304-3975(97)00190-4.

\bibitem{kuriyamaCFG}
T.~Kuriyama.
\newblock Distributional learning of context-free languages under fixed
finite-monoid typing.
\newblock Revised manuscript prepared for \emph{Theoretical Computer Science},
2026; public preprint arXiv:1409.6247v4 [cs.FL].
\newblock The revised proof of Proposition~6.3 and Appendix~C is
available in \href{https://github.com/growupkuriyama-hub/Papers/blob/7f4c680e3a896b41b748970d75d5084fa6e08e6d/01_fixed-h-cfg/main.tex}{repository revision 7f4c680e}.

\bibitem{McNaughtonPapert1971}
R.~McNaughton and S.~Papert.
\newblock \emph{Counter-Free Automata}.
\newblock MIT Press, Cambridge, MA, 1971.

\bibitem{Pin1986}
J.-E. Pin.
\newblock {\em Varieties of Formal Languages}.
\newblock North Oxford Academic, London, and Plenum, New York, 1986.

\bibitem{SekiEtAl1991}
H.~Seki, T.~Matsumura, M.~Fujii, and T.~Kasami.
\newblock On multiple context-free grammars.
\newblock {\em Theoretical Computer Science}, 88(2):191--229, 1991.
\newblock doi:10.1016/0304-3975(91)90374-B.

\bibitem{Kallmeyer2010}
L.~Kallmeyer.
\newblock {\em Parsing Beyond Context-Free Grammars}.
\newblock Springer, Berlin, Heidelberg, 2010.
\newblock doi:10.1007/978-3-642-14846-0.

\bibitem{Kanazawa2019Ogden}
M.~Kanazawa.
\newblock Ogden's lemma, multiple context-free grammars, and the control language hierarchy.
\newblock {\em Information and Computation}, 269:104449, 2019.
\newblock doi:10.1016/j.ic.2019.104449.

\bibitem{GomezRodriguezEtAl2009}
C.~G\'omez-Rodr\'iguez, M.~Kuhlmann, G.~Satta, and D.~Weir.
\newblock Optimal reduction of rule length in linear context-free rewriting systems.
\newblock In {\em Proceedings of Human Language Technologies: The 2009 Annual
Conference of the North American Chapter of the Association for Computational
Linguistics}, pp.~539--547. Association for Computational Linguistics, 2009.


\end{thebibliography}
\normalsize

\end{document}